\documentclass[aps,preprint]{revtex4}%
\usepackage{amsfonts}
\usepackage{amsmath}
\usepackage{amssymb}
\usepackage{graphicx}%
\setcounter{MaxMatrixCols}{30}
%TCIDATA{OutputFilter=latex2.dll}
%TCIDATA{Version=5.50.0.2953}
%TCIDATA{CSTFile=revtex4.cst}
%TCIDATA{Created=Thursday, June 26, 2008 11:45:27}
%TCIDATA{LastRevised=Friday, November 13, 2009 15:22:20}
%TCIDATA{<META NAME="GraphicsSave" CONTENT="32">}
%TCIDATA{<META NAME="SaveForMode" CONTENT="1">}
%TCIDATA{BibliographyScheme=Manual}
%TCIDATA{<META NAME="DocumentShell" CONTENT="Articles\SW\REVTeX 4">}
%BeginMSIPreambleData
\providecommand{\U}[1]{\protect\rule{.1in}{.1in}}
%EndMSIPreambleData
\newtheorem{theorem}{Theorem}
\newtheorem{acknowledgement}[theorem]{Acknowledgement}
\newtheorem{algorithm}[theorem]{Algorithm}
\newtheorem{axiom}[theorem]{Axiom}
\newtheorem{claim}[theorem]{Claim}
\newtheorem{conclusion}[theorem]{Conclusion}
\newtheorem{condition}[theorem]{Condition}
\newtheorem{conjecture}[theorem]{Conjecture}
\newtheorem{corollary}[theorem]{Corollary}
\newtheorem{criterion}[theorem]{Criterion}
\newtheorem{definition}[theorem]{Definition}
\newtheorem{example}[theorem]{Example}
\newtheorem{exercise}[theorem]{Exercise}
\newtheorem{lemma}[theorem]{Lemma}
\newtheorem{notation}[theorem]{Notation}
\newtheorem{problem}[theorem]{Problem}
\newtheorem{proposition}[theorem]{Proposition}
\newtheorem{remark}[theorem]{Remark}
\newtheorem{solution}[theorem]{Solution}
\newtheorem{summary}[theorem]{Summary}
\newenvironment{proof}[1][Proof]{\noindent\textbf{#1.} }{\ \rule{0.5em}{0.5em}}
\begin{document}
\preprint{ }
\title[Short title for running header]{Tangent and Normal Spaces}
\author{Brian Weiner}
\affiliation{Penn State}
\startpage{1}
\maketitle
\tableofcontents


\subsection{General Setting}

Consider the, in general nonlinear, map
\begin{equation}
\boldsymbol{g}:\mathbb{R}^{n}\longrightarrow\mathbb{R}^{m}%
\end{equation}
given by%
\begin{equation}
\boldsymbol{g}\left(  \boldsymbol{A}\right)  =\left(
\begin{array}
[c]{c}%
g_{1}\left(  \boldsymbol{A}\right) \\
\vdots\\
g_{m}\left(  \boldsymbol{A}\right)
\end{array}
\right)  =\left(
\begin{array}
[c]{c}%
\rho_{1}\\
\vdots\\
\rho_{m}%
\end{array}
\right)  =\boldsymbol{\rho}\in\mathbb{R}^{m};\boldsymbol{A\in}\mathbb{R}^{n},
\end{equation}
that defines a manifold
\begin{equation}
\mathfrak{M}=\left\{  \left.  \boldsymbol{A}\right\vert \boldsymbol{g}\left(
\boldsymbol{A}\right)  =\boldsymbol{\rho}\right\}  ,
\end{equation}
which is a level surface in $\mathbb{R}^{n+m}.$ In this case both the domain
and range are inner product spaces.

\subsection{Density Functional Example}

As a special case one can consider the bilinear function
\begin{equation}
\boldsymbol{f}\left(  A\right)  =\left(  \left.  A\right\vert \widehat
{\Phi^{\dagger}\Phi_{\varphi}}A\right)  ,\quad A\in\mathcal{B}_{sa}\left(
\mathcal{H}\right)  \cong\mathbb{R}^{n^{2}};n=\binom{m}{N}^{2};\quad
\boldsymbol{f}\left(  A\right)  \in\mathbb{R}^{m}%
\end{equation}
where $\Phi^{\dagger}\Phi_{\varphi}$ is a fixed vector operator with
components
\begin{equation}
\Phi^{\dagger}\Phi_{\varphi j}=\Phi\left(  \varphi_{j}\right)  ^{\dagger}%
\Phi\left(  \varphi_{j}\right)
\end{equation}
formed by the conb $\left\{  \left.  \varphi_{j}\right\vert 1\leq j\leq
m\right\}  $ for $\mathcal{H}$ and
\begin{equation}
\widehat{\Phi^{\dagger}\Phi_{\varphi}}_{j}\left\vert A\right)  =\frac{1}%
{2}\left\vert \left[  \Phi\left(  \varphi_{j}\right)  ^{\dagger}\Phi\left(
\varphi_{j}\right)  ,A\right]  _{+}\right)  \equiv\widehat{\Phi^{\dagger}%
\Phi_{\varphi}}_{j}A=\frac{1}{2}\left[  \Phi\left(  \varphi_{j}\right)
^{\dagger}\Phi\left(  \varphi_{j}\right)  ,A\right]  _{+}.
\end{equation}
As $m\rightarrow\infty$ the operators $\Phi\left(  \varphi_{j}\right)
^{\dagger}\Phi\left(  \varphi_{j}\right)  \rightarrow\Phi\left(
\delta_{\mathbf{r}}\right)  ^{\dagger}\Phi\left(  \delta_{\mathbf{r}}\right)
$ and $\Phi^{\dagger}\Phi_{\varphi}\rightarrow\Phi^{\dagger}\Phi_{\delta}$
i.e. the vector operator tends to a distributional operator%
\begin{align}
\Phi^{\dagger}\Phi_{\varphi}  &  :\mathbb{N}\rightarrow\mathcal{B}\left(
\mathcal{H}\right) \\
\Phi^{\dagger}\Phi_{\delta}  &  :\mathcal{D}_{\delta}\equiv\mathbb{R}%
^{3}\rightarrow\mathcal{B}\left(  \mathcal{H}\right)
\end{align}
and $\Phi^{\dagger}\Phi_{\varphi}$ defines the super operator $\widehat
{\Phi^{\dagger}\Phi_{\varphi}}$ by symmetrized multiplication. Noting that
\begin{equation}
A=\sum_{1\leq p,q\leq n}A_{pq}\left\vert \Psi_{p}\right\rangle \left\langle
\Psi_{q}\right\vert
\end{equation}
this quadratic function can written in detail as%
\begin{align}
&  \boldsymbol{f}_{j}\left(  A\right)  =\left(  A\left\vert \widehat
{\Phi^{\dagger}\Phi_{\varphi j}}A\right.  \right)  =\frac{1}{2}%
\operatorname{Tr}\left\{  A\left[  \Phi\left(  \varphi_{j}\right)  ^{\dagger
}\Phi\left(  \varphi_{j}\right)  ,A\right]  _{+}\right\} \\
&  =\frac{1}{2}\operatorname{Tr}\left\{  A\left(  \Phi\left(  \varphi
_{j}\right)  ^{\dagger}\Phi\left(  \varphi_{j}\right)  A+A\Phi\left(
\varphi_{j}\right)  ^{\dagger}\Phi\left(  \varphi_{j}\right)  \right)
\right\}  =\operatorname{Tr}\left\{  \Phi\left(  \varphi_{j}\right)
^{\dagger}\Phi\left(  \varphi_{j}\right)  A^{2}\right\} \\
&  =\sum_{1\leq p,q,s\leq n}\left\langle \Psi_{s}\left\vert \Phi\left(
\varphi_{j}\right)  ^{\dagger}\Phi\left(  \varphi_{j}\right)  \Psi_{p}\right.
\right\rangle A_{pq}A_{qs}%
\end{align}
The first derrivatives of $\boldsymbol{f}$ at the point $A$ are then
\begin{align}
\frac{\partial\boldsymbol{f}_{j}\left(  A\right)  }{\partial A_{\alpha\beta}}
&  =\sum_{1\leq p,q,s\leq n}\left\langle \Psi_{s}\left\vert \Phi\left(
\varphi_{j}\right)  ^{\dagger}\Phi\left(  \varphi_{j}\right)  \Psi_{p}\right.
\right\rangle \left\{  \delta_{\alpha p}\delta_{\beta q}A_{qs}+A_{pq}%
\delta_{\alpha q}\delta_{\beta s}\right\} \\
&  =\sum_{1\leq s\leq n}\left\{  A_{\beta s}\left\langle \Psi_{s}\left\vert
\Phi\left(  \varphi_{j}\right)  ^{\dagger}\Phi\left(  \varphi_{j}\right)
\Psi_{\alpha}\right.  \right\rangle +\left\langle \Psi_{\beta}\left\vert
\Phi\left(  \varphi_{j}\right)  ^{\dagger}\Phi\left(  \varphi_{j}\right)
\Psi_{s}\right.  \right\rangle A_{s\alpha}\right\} \\
&  =\sum_{1\leq s\leq n}\left\{  A_{\beta s}\left\langle \Psi_{s}\left\vert
\Phi\left(  \varphi_{j}\right)  ^{\dagger}\Phi\left(  \varphi_{j}\right)
\Psi_{\alpha}\right.  \right\rangle +\left\langle \Psi_{\beta}\left\vert
\Phi\left(  \varphi_{j}\right)  ^{\dagger}\Phi\left(  \varphi_{j}\right)
\Psi_{s}\right.  \right\rangle A_{s\alpha}\right\}
\end{align}
and the second
\begin{equation}
\frac{\partial^{2}\boldsymbol{f}_{j}\left(  A\right)  }{\partial A_{\lambda
\mu}\partial A_{\alpha\beta}}=\left\langle \Psi_{\mu}\left\vert \Phi\left(
\varphi_{j}\right)  ^{\dagger}\Phi\left(  \varphi_{j}\right)  \Psi_{\alpha
}\right.  \right\rangle \delta_{\lambda\beta}+\left\langle \Psi_{\beta
}\left\vert \Phi\left(  \varphi_{j}\right)  ^{\dagger}\Phi\left(  \varphi
_{j}\right)  \Psi_{\lambda}\right.  \right\rangle \delta_{\mu\alpha}.
\end{equation}
For
\begin{equation}
A\rightarrow A+\delta A
\end{equation}
(if $\delta A$ belongs to the tangent space then $A+\delta A$ does not belong
to the surface), i.e. expanding for each component in a small nbhd of $A,$ we
get
\begin{align}
&  \boldsymbol{f}_{j}\left(  A+\delta A\right)  =\sum_{1\leq p,q,s\leq
n}\left\langle \Psi_{s}\left\vert \Phi\left(  \varphi_{j}\right)  ^{\dagger
}\Phi\left(  \varphi_{j}\right)  \Psi_{p}\right.  \right\rangle \left(
A+\delta A\right)  _{pq}\left(  A+\delta A\right)  _{qs}\\
&  =\sum_{1\leq p,q,s\leq n}\left\langle \Psi_{s}\left\vert \Phi\left(
\varphi_{j}\right)  ^{\dagger}\Phi\left(  \varphi_{j}\right)  \Psi_{p}\right.
\right\rangle A_{pq}A_{qs}\\
&  +\sum_{1\leq p,q,s\leq n}\left\langle \Psi_{s}\left\vert \Phi\left(
\varphi_{j}\right)  ^{\dagger}\Phi\left(  \varphi_{j}\right)  \Psi_{p}\right.
\right\rangle \left\{  \delta A_{pq}A_{qs}+A_{pq}\delta A_{qs}\right\} \\
&  +\sum_{1\leq p,q,s\leq n}\left\langle \Psi_{s}\left\vert \Phi\left(
\varphi_{j}\right)  ^{\dagger}\Phi\left(  \varphi_{j}\right)  \Psi_{p}\right.
\right\rangle \delta A_{pq}\delta A_{qs}%
\end{align}
giving%
\begin{align}
&  =\rho_{j}+\sum_{1\leq p,q,s\leq n}\left\{  A_{qs}\left\langle \Psi
_{s}\left\vert \Phi\left(  \varphi_{j}\right)  ^{\dagger}\Phi\left(
\varphi_{j}\right)  \Psi_{p}\right.  \right\rangle \delta A_{pq}+\delta
A_{qs}\left\langle \Psi_{s}\left\vert \Phi\left(  \varphi_{j}\right)
^{\dagger}\Phi\left(  \varphi_{j}\right)  \Psi_{p}\right.  \right\rangle
A_{pq}\right\} \\
&  +\sum_{1\leq p,q,s\leq n}\delta A_{qs}\left\langle \Psi_{s}\left\vert
\Phi\left(  \varphi_{j}\right)  ^{\dagger}\Phi\left(  \varphi_{j}\right)
\Psi_{p}\right.  \right\rangle \delta A_{pq}%
\end{align}
$\left(  \text{the last term can be written as }\sum_{1\leq p,s\leq
n}\left\langle \Psi_{s}\left\vert \Phi\left(  \varphi_{j}\right)  ^{\dagger
}\Phi\left(  \varphi_{j}\right)  \Psi_{p}\right.  \right\rangle \left(  \delta
A\right)  _{ps}^{2}\right)  $ producing
\begin{align}
\boldsymbol{f}\left(  A\right)   &  =\boldsymbol{\rho}+2\left(  j\updownarrow%
\begin{array}
[c]{ccc}%
\vdots & \cdots & \vdots\\
\vdots & \left\langle \Psi_{q}\left\vert A\Phi\left(  \varphi_{j}\right)
^{\dagger}\Phi\left(  \varphi_{j}\right)  \Psi_{p}\right.  \right\rangle  &
\vdots\\
\vdots & \cdots & \vdots
\end{array}
\right)  \left(
\begin{array}
[c]{c}%
\vdots\\
\delta A_{pq}\\
\vdots
\end{array}
\right) \\
&  +\frac{1}{2}\left(
\begin{array}
[c]{ccc}%
\cdots & \delta A_{pq} & \cdots
\end{array}
\right)  \left(
\begin{array}
[c]{ccc}%
\vdots & \cdots & \vdots\\
\vdots & \left\langle \Psi_{s}\left\vert \Phi\left(  \varphi_{j}\right)
^{\dagger}\Phi\left(  \varphi_{j}\right)  \Psi_{p}\right.  \right\rangle
\delta_{rq} & \vdots\\
\vdots & \cdots & \vdots
\end{array}
\right)  \left(
\begin{array}
[c]{c}%
\vdots\\
\delta A_{rs}\\
\vdots
\end{array}
\right)
\end{align}


\subsection{General Function}

Assuming the function $\boldsymbol{g}\left(  \boldsymbol{A}\right)  $ is
Fr\'{e}chet differentiable
\begin{align}
\left[  d\boldsymbol{g}\left(  \boldsymbol{A}\right)  \right]  \left(
\delta\boldsymbol{A}\right)   &  =\left(
\begin{array}
[c]{ccc}%
\frac{\partial g_{1}\left(  \boldsymbol{A}\right)  }{\partial A_{1}} & \cdots
& \frac{\partial g_{1}\left(  \boldsymbol{A}\right)  }{\partial A_{n}}\\
\vdots & \ddots & \vdots\\
\frac{\partial g_{m}\left(  \boldsymbol{A}\right)  }{\partial A_{1}} & \cdots
& \frac{\partial g_{m}\left(  \boldsymbol{A}\right)  }{\partial A_{n}}%
\end{array}
\right)  \left(
\begin{array}
[c]{c}%
\delta A_{1}\\
\vdots\\
\delta A_{n}%
\end{array}
\right)  =\mathbb{D}\boldsymbol{g}\left(  \boldsymbol{A}\right)  \left\vert
\delta\boldsymbol{A}\right) \nonumber\\
&  =\left(
\begin{array}
[c]{c}%
\nabla g_{1}\left(  \boldsymbol{A}\right) \\
\vdots\\
\nabla g_{m}\left(  \boldsymbol{A}\right)
\end{array}
\right)  \left(
\begin{array}
[c]{c}%
\delta A_{1}\\
\vdots\\
\delta A_{n}%
\end{array}
\right)  =\left(
\begin{array}
[c]{c}%
\delta\rho_{1}\\
\vdots\\
\delta\rho_{m}%
\end{array}
\right)  \in\mathbb{R}^{m}.
\end{align}%
\begin{align}
&  \left[  d^{2}\boldsymbol{g}\left(  \boldsymbol{A}\right)  \right]  \left(
\delta\boldsymbol{A}\right)  =\frac{1}{2}\left(
\begin{array}
[c]{ccc}%
\delta A_{1} & \cdots & \delta A_{n}%
\end{array}
\right)  \left(
\begin{array}
[c]{c}%
\nabla^{2}g_{1}\left(  \boldsymbol{A}\right) \\
\vdots\\
\nabla^{2}g_{m}\left(  \boldsymbol{A}\right)
\end{array}
\right)  \left(
\begin{array}
[c]{c}%
\delta A_{1}\\
\vdots\\
\delta A_{n}%
\end{array}
\right) \nonumber\\
&  =\frac{1}{2}\left(
\begin{array}
[c]{c}%
\left(
\begin{array}
[c]{ccc}%
\delta A_{1} & \cdots & \delta A_{n}%
\end{array}
\right)  \nabla^{2}g_{1}\left(  \boldsymbol{A}\right)  \left(
\begin{array}
[c]{c}%
\delta A_{1}\\
\vdots\\
\delta A_{n}%
\end{array}
\right) \\
\vdots\\
\left(
\begin{array}
[c]{ccc}%
\delta A_{1} & \cdots & \delta A_{n}%
\end{array}
\right)  \nabla^{2}g_{m}\left(  \boldsymbol{A}\right)  \left(
\begin{array}
[c]{c}%
\delta A_{1}\\
\vdots\\
\delta A_{n}%
\end{array}
\right)
\end{array}
\right)  =\left(
\begin{array}
[c]{c}%
\delta^{2}\rho_{1}\\
\vdots\\
\delta^{2}\rho_{m}%
\end{array}
\right)  \in\mathbb{R}^{m}.
\end{align}


\subsection{Density Functional Example}

For the above example
\begin{align}
\left(  \nabla f_{j}\left(  A\right)  \right)  _{pq}  &  =\frac{\partial
f_{j}\left(  A\right)  }{\partial A_{pq}}=\left(  \left.  A\right\vert
\Phi\left(  \varphi_{j}\right)  ^{\dagger}\Phi\left(  \varphi_{j}\right)
\left\vert \varphi_{p}\right\rangle \left\langle \varphi_{q}\right\vert
\right)  =\operatorname{Tr}\left\{  A\Phi\left(  \varphi_{j}\right)
^{\dagger}\Phi\left(  \varphi_{j}\right)  \left\vert \varphi_{p}\right\rangle
\left\langle \varphi_{q}\right\vert \right\} \\
&  =\left\langle \varphi_{q}\left\vert A\Phi\left(  \varphi_{j}\right)
^{\dagger}\Phi\left(  \varphi_{j}\right)  \varphi_{p}\right.  \right\rangle
;\quad1\leq j\leq m;\,1\leq p,q\leq\sqrt{n}%
\end{align}
and
\begin{equation}
d\boldsymbol{f}\left(  A\right)  =\left(  \widehat{\Phi^{\dagger}\Phi
_{\varphi}}A\right\vert =\frac{1}{2}\left(  \left[  \Phi^{\dagger}%
\Phi_{\varphi},A\right]  _{+}\right\vert .
\end{equation}
The second derrivatives are%
\begin{equation}
\left(  \nabla^{2}f_{j}\left(  A\right)  \right)  _{pqrs}=\left\langle
\Psi_{s}\left\vert \Phi\left(  \varphi_{j}\right)  ^{\dagger}\Phi\left(
\varphi_{j}\right)  \Psi_{p}\right.  \right\rangle \delta_{rq};\quad1\leq
p,q,r,s\leq\sqrt{n}.
\end{equation}


\subsection{General Function}

The \emph{kernel}, $\ker\left\{  d\boldsymbol{g}\left(  \boldsymbol{A}\right)
\right\}  ,$of $d\boldsymbol{g}\left(  \boldsymbol{A}\right)  $ is defined as%
\begin{equation}
\ker\left\{  d\boldsymbol{g}\left(  \boldsymbol{A}\right)  \right\}  =\left\{
\left.  \delta\boldsymbol{A}\right\vert \left[  d\boldsymbol{g}\left(
\boldsymbol{A}\right)  \right]  \left(  \delta\boldsymbol{A}\right)
=\boldsymbol{0}\right\}  \subseteq\mathbb{R}^{n}.
\end{equation}
The \emph{tangent space} $\mathcal{T}_{\boldsymbol{A}}$ of $\boldsymbol{g}$ at
the point $\boldsymbol{A}$ is contained in this kernel. A tangent vector
$\delta\boldsymbol{A}$ at the point $\boldsymbol{A}$ is defined by considering
curves $\in\mathfrak{M}$%
\begin{align}
\boldsymbol{g}\left(  \boldsymbol{x}\left(  t\right)  \right)   &
=\mathbf{\rho}\,,t\in\left[  -\delta,\delta\right] \nonumber\\
\boldsymbol{x}\left(  0\right)   &  =\boldsymbol{A}\nonumber\\
\boldsymbol{x}  &  :\left[  -\delta,\delta\right]  \longrightarrow
\mathbb{R}^{n}%
\end{align}
and is defined to be formed by the equivalence class of velocity vectors to
such curves at the point $\boldsymbol{A}$ i.e. $\left\{  \frac{d\boldsymbol{x}%
\left(  0\right)  }{dt}\right\}  .$ The Taylor series of $\boldsymbol{x}$
around $0$ is
\begin{equation}
\boldsymbol{x}\left(  t\right)  =\boldsymbol{A}+\frac{d\boldsymbol{x}\left(
0\right)  }{dt}t+\frac{1}{2}\frac{d^{2}\boldsymbol{x}\left(  0\right)
}{dt^{2}}t^{2}+\cdots
\end{equation}
and of $\boldsymbol{g}$ restricted to the curve $\boldsymbol{x}$ that lies in
the surface $\mathfrak{M}$ is
\begin{equation}
\boldsymbol{g}\left(  \boldsymbol{x}\left(  t\right)  \right)
=\boldsymbol{\rho}+\frac{d\boldsymbol{g}\left(  \boldsymbol{x}\left(
0\right)  \right)  }{dt}t+\frac{1}{2}\frac{d^{2}\boldsymbol{g}\left(
\boldsymbol{x}\left(  0\right)  \right)  }{dt^{2}}t^{2}+\cdots
=\boldsymbol{\rho}%
\end{equation}
where%
\begin{align}
\frac{d\boldsymbol{g}\left(  \boldsymbol{x}\left(  0\right)  \right)  }{dt}
&  =\mathbb{D}_{x}\boldsymbol{g}\left(  \boldsymbol{x}\left(  0\right)
\right)  \frac{d\boldsymbol{x}\left(  0\right)  }{dt}\nonumber\\
\frac{d^{2}\boldsymbol{g}\left(  \boldsymbol{x}\left(  0\right)  \right)
}{dt^{2}}  &  =\left.  \left\{  \frac{d}{dt}\mathbb{D}_{x}\boldsymbol{g}%
\left(  \boldsymbol{x}\left(  t\right)  \right)  \right\}  \right\vert
_{t=0}\frac{d\boldsymbol{x}\left(  0\right)  }{dt}+\mathbb{D}_{x}%
\boldsymbol{g}\left(  \boldsymbol{x}\left(  t\right)  \right)  \frac
{d^{2}\boldsymbol{x}\left(  0\right)  }{dt^{2}}\nonumber\\
&  =\left(  \frac{d\boldsymbol{x}\left(  0\right)  }{dt}\right)
^{t}\mathbb{D}_{x}^{2}\boldsymbol{g}\left(  \boldsymbol{x}\left(  0\right)
\right)  \frac{d\boldsymbol{x}\left(  0\right)  }{dt}+\mathbb{D}%
_{x}\boldsymbol{g}\left(  \boldsymbol{x}\left(  0\right)  \right)  \frac
{d^{2}\boldsymbol{x}\left(  0\right)  }{dt^{2}}.
\end{align}
Hence
\begin{equation}
\frac{d\boldsymbol{g}\left(  \boldsymbol{x}\left(  t\right)  \right)  }%
{dt}=\frac{d\boldsymbol{g}\left(  \boldsymbol{x}\left(  0\right)  \right)
}{dt}+\frac{d^{2}\boldsymbol{g}\left(  \boldsymbol{x}\left(  0\right)
\right)  }{dt^{2}}t+\cdots=\mathbf{0}%
\end{equation}
and \emph{in particular }at $t=0$\emph{ }%
\begin{equation}
\frac{d\boldsymbol{g}\left(  \boldsymbol{x}\left(  0\right)  \right)  }%
{dt}=\mathbb{D}_{x}\boldsymbol{g}\left(  \boldsymbol{x}\left(  0\right)
\right)  \frac{d\boldsymbol{x}\left(  0\right)  }{dt}=\mathbf{0},
\end{equation}
i.e.
\begin{equation}
\frac{d\boldsymbol{x}\left(  0\right)  }{dt}\in\ker\left\{  \mathbb{D}%
_{x}\boldsymbol{g}\left(  \boldsymbol{x}\left(  0\right)  \right)  \right\}  .
\end{equation}
The derrivatives wrt $x$ are the same as wrt $\boldsymbol{A}$ thus
\begin{align}
\mathbb{D}_{x}\boldsymbol{g}\left(  \boldsymbol{x}\left(  0\right)  \right)
&  =\mathbb{D}\boldsymbol{g}\left(  \boldsymbol{A}\right) \nonumber\\
\mathbb{D}_{x}^{2}\boldsymbol{g}\left(  \boldsymbol{x}\left(  0\right)
\right)   &  =\mathbb{D}^{2}\boldsymbol{g}\left(  \boldsymbol{A}\right)  ,
\end{align}
thus
\begin{equation}
\frac{d\boldsymbol{x}\left(  0\right)  }{dt}\in\ker\left\{  \mathbb{D}%
\boldsymbol{g}\left(  \boldsymbol{A}\right)  \right\}  .
\end{equation}
One should note that
\begin{equation}
\delta\boldsymbol{A}\in\ker\left\{  \mathbb{D}\boldsymbol{g}\left(
\boldsymbol{A}\right)  \right\}
\end{equation}
\emph{does not imply} that $\exists$ $\frac{d\boldsymbol{x}\left(  0\right)
}{dt}$ $\backepsilon$ $\delta\boldsymbol{A}=\frac{d\boldsymbol{x}\left(
0\right)  }{dt}$ i.e. this condition is necessary but not sufficient for
$\delta\boldsymbol{A}$ to be a tangent vector. One can write the above as
\begin{equation}
\left[  d\boldsymbol{g}\left(  \boldsymbol{x}\left(  0\right)  \right)
\right]  \left(  \delta\boldsymbol{A}\right)  =\mathbb{D}_{x}\boldsymbol{g}%
\left(  \boldsymbol{x}\left(  0\right)  \right)  \left(
\begin{array}
[c]{c}%
\frac{\boldsymbol{d}x_{1}\left(  0\right)  }{\boldsymbol{d}t}\\
\vdots\\
\frac{\boldsymbol{d}x_{n}\left(  0\right)  }{\boldsymbol{d}t}%
\end{array}
\right)  =\boldsymbol{0}.
\end{equation}
In detail this gives
\begin{align}
\mathbb{D}_{x}\boldsymbol{g}\left(  \boldsymbol{x}\left(  0\right)  \right)
\left(
\begin{array}
[c]{c}%
\frac{\boldsymbol{d}x_{1}\left(  0\right)  }{\boldsymbol{d}t}\\
\vdots\\
\frac{\boldsymbol{d}x_{n}\left(  0\right)  }{\boldsymbol{d}t}%
\end{array}
\right)   &  =\left(
\begin{array}
[c]{c}%
\nabla g_{1}\left(  \boldsymbol{A}\right) \\
\vdots\\
\nabla g_{m}\left(  \boldsymbol{A}\right)
\end{array}
\right)  \left(
\begin{array}
[c]{c}%
\frac{\boldsymbol{d}x_{1}\left(  0\right)  }{\boldsymbol{d}t}\\
\vdots\\
\frac{\boldsymbol{d}x_{n}\left(  0\right)  }{\boldsymbol{d}t}%
\end{array}
\right) \nonumber\\
&  =\left(
\begin{array}
[c]{ccc}%
\frac{\partial g_{1}\left(  \boldsymbol{A}\right)  }{\partial\boldsymbol{A}%
_{1}} & \cdots & \frac{\partial g_{1}\left(  \boldsymbol{A}\right)  }%
{\partial\boldsymbol{A}_{n}}\\
\vdots & \ddots & \vdots\\
\frac{\partial g_{m}\left(  \boldsymbol{A}\right)  }{\partial\boldsymbol{A}%
_{1}} & \cdots & \frac{\partial g_{m}\left(  \boldsymbol{A}\right)  }%
{\partial\boldsymbol{A}_{n}}%
\end{array}
\right)  \left(
\begin{array}
[c]{c}%
\frac{\boldsymbol{d}x_{1}\left(  0\right)  }{\boldsymbol{d}t}\\
\vdots\\
\frac{\boldsymbol{d}x_{n}\left(  0\right)  }{\boldsymbol{d}t}%
\end{array}
\right)  =0.
\end{align}


\subsection{Density Functional Example}

For the example $\boldsymbol{f}\left(  A\right)  $ this gives that if $\delta
A\in\mathcal{T}_{A}\mathfrak{\ }$then
\begin{equation}
\left[  d\boldsymbol{f}\left(  A\right)  \right]  \left(  \delta A\right)
=\mathbf{0}, \label{Ker}%
\end{equation}
i.e. if $\delta A\in\mathcal{T}_{A}\mathfrak{\ }$then $\left\vert \delta
A\right)  \bot\left\vert \left[  \Phi\left(  \varphi_{j}\right)  ^{\dagger
}\Phi\left(  \varphi_{j}\right)  ,A\right]  _{+}\right)  \forall j,$but being
orthogonal $\forall j$ does not mean it is in the tangent space, i.e. does not
ensure that $\delta A=\dot{x}\left(  0\right)  $, which can thus be smaller
than the kernel.

One should note that $\left(  \widehat{\Phi^{\dagger}\Phi_{\varphi}%
}A\right\vert $ is not an operator but a \emph{vector} \emph{operator }with
operator components $\left[  \Phi\left(  \varphi_{j}\right)  ^{\dagger}%
\Phi\left(  \varphi_{j}\right)  ,A\right]  _{+},$ $1\leq j\leq m$, so the
condition eq. $\left(  \ref{Ker}\right)  $ gives $m$ equations%
\begin{equation}
\left(  A\left\vert \widehat{\Phi\left(  \varphi_{j}\right)  ^{\dagger}%
\Phi\left(  \varphi_{j}\right)  }\delta A\right.  \right)  =\operatorname{Tr}%
\left\{  A\Phi\left(  \varphi_{j}\right)  ^{\dagger}\Phi\left(  \varphi
_{j}\right)  \delta A\right\}  =0;\quad1\leq j\leq m.
\end{equation}
Note that if $\delta A\in\mathcal{T}_{A}$ then $A+\delta A$ does not belong to
the surface.

\subsubsection{Tangent Vectors in the DFT\ example}

The curves in the surface satisfy%

\begin{equation}
\left(  x\left(  t\right)  \left\vert \widehat{\Phi\left(  \varphi_{j}\right)
^{\dagger}\Phi\left(  \varphi_{j}\right)  }x\left(  t\right)  \right.
\right)  =\operatorname{Tr}\left\{  x\left(  t\right)  \Phi\left(  \varphi
_{j}\right)  ^{\dagger}\Phi\left(  \varphi_{j}\right)  x\left(  t\right)
\right\}  =\rho_{j};\quad1\leq j\leq m
\end{equation}
and belong to the surface $\mathfrak{M\subset}\mathcal{B}_{sa}\left(
\mathcal{H}^{N}\right)  $
\begin{equation}
\mathfrak{M}=\left\{  A\left\vert \left(  A\left\vert \widehat{\Phi\left(
\varphi_{j}\right)  ^{\dagger}\Phi\left(  \varphi_{j}\right)  }A\right.
\right)  =\operatorname{Tr}\left\{  \Phi\left(  \varphi_{j}\right)  ^{\dagger
}\Phi\left(  \varphi_{j}\right)  A^{2}\right\}  =\rho_{j};\quad1\leq j\leq
m\right.  \right\}  .
\end{equation}
Note that
\begin{equation}
\left(  A\left\vert \widehat{\Phi\left(  \varphi_{j}\right)  ^{\dagger}%
\Phi\left(  \varphi_{j}\right)  }A\right.  \right)  =\left\langle A\left\vert
\Phi\left(  \varphi_{j}\right)  ^{\dagger}\Phi\left(  \varphi_{j}\right)
A\right.  \right\rangle _{HS}%
\end{equation}
in terms of the Hilbert-Schmidt product
\begin{equation}
\left\langle \cdot\left\vert \cdot\right.  \right\rangle _{HS}:\mathcal{B}%
_{2}\left(  \mathcal{H}^{N}\right)  \times\mathcal{B}_{2}\left(
\mathcal{H}^{N}\right)  \rightarrow\mathbb{C}.
\end{equation}
The real symmetric inner product
\begin{equation}
\left(  \cdot\left\vert \cdot\right.  \right)  :\mathcal{B}_{2sa}\left(
\mathcal{H}^{N}\right)  \times\mathcal{B}_{2sa}\left(  \mathcal{H}^{N}\right)
\rightarrow\mathbb{R}%
\end{equation}
has the invariance group $O\left(  \mathcal{B}_{2sa}\left(  \mathcal{H}%
^{N}\right)  \right)  \equiv O\left(  n,\mathbb{R}\right)  .$The group of
inner automorphisms%
\begin{align}
\widehat{U}  &  :\mathcal{B}_{2sa}\left(  \mathcal{H}^{N}\right)
\rightarrow\mathcal{B}_{2sa}\left(  \mathcal{H}^{N}\right) \\
\widehat{U}\left(  A\right)   &  =U^{\dagger}AU;\,U\in U\left(  \mathcal{H}%
^{N}\right)  ;\,A\in\mathcal{B}_{2sa}\left(  \mathcal{H}^{N}\right)
\end{align}
is contained in $O\left(  \mathcal{B}_{2sa}\left(  \mathcal{H}^{N}\right)
\right)  .$ The operators
\begin{equation}
\left\{  \left.  \widehat{\Phi\left(  \varphi_{j}\right)  ^{\dagger}%
\Phi\left(  \varphi_{j}\right)  }\right\vert 1\leq j\leq m\right\}
\end{equation}
form a commuting set thus are simultaneously diagonalizable. The conditions

%

\begin{align}
&  \left(  A\left\vert \widehat{O}\left(  t\right)  ^{t}\widehat{\Phi\left(
\varphi_{j}\right)  ^{\dagger}\Phi\left(  \varphi_{j}\right)  }\widehat
{O}\left(  t\right)  A\right.  \right)  =\left(  A\left\vert \widehat
{\Phi\left(  \varphi_{j}\right)  ^{\dagger}\Phi\left(  \varphi_{j}\right)
}A\right.  \right)  \equiv\label{invar}\\
&  \operatorname{Tr}\left\{  \left(  \widehat{O}\left(  t\right)  \left(
A\right)  \right)  \Phi\left(  \varphi_{j}\right)  ^{\dagger}\Phi\left(
\varphi_{j}\right)  \left(  \widehat{O}\left(  t\right)  \left(  A\right)
\right)  \right\}  =\operatorname{Tr}\left\{  A\Phi\left(  \varphi_{j}\right)
^{\dagger}\Phi\left(  \varphi_{j}\right)  A\right\}
\end{align}%
\begin{equation}
\qquad\qquad\qquad\qquad\qquad\qquad\qquad\qquad\qquad\qquad\qquad1\leq j\leq
m,\widehat{O}\left(  t\right)  \in\mathfrak{G,}%
\end{equation}
where
\begin{equation}
\widehat{O}\left(  t\right)  =\exp\left\{  \sum_{1\leq k\leq\nu}\lambda
_{k}\widehat{\Lambda}_{k}\right\}  ;\widehat{\Lambda}_{k}=-\widehat{\Lambda
}_{k}^{t}\,,\lambda_{k}\in\mathbb{R},1\leq k\leq\nu
\end{equation}
and $\left\{  \left.  \widehat{\Lambda}_{k}\right\vert 1\leq k\leq\nu\right\}
$ are the generators of the Lie algebra, $\mathcal{G},$ of $\mathfrak{G}$. The
conditions eq $\left(  \ref{invar}\right)  $ imply the necessary but not
sufficient conditions
\begin{equation}
\left(  A\left\vert \left[  \widehat{\Phi\left(  \varphi_{j}\right)
^{\dagger}\Phi\left(  \varphi_{j}\right)  },\widehat{\Lambda}_{k}\right]
\right.  A\right)  =0;1\leq k\leq\nu\,\text{, }1\leq j\leq m,
\end{equation}
which can be written as
\begin{align}
0  &  =\operatorname{Tr}\left\{  \left[  \widehat{\Phi\left(  \varphi
_{j}\right)  ^{\dagger}\Phi\left(  \varphi_{j}\right)  },\widehat{\Lambda}%
_{k}\right]  \left\vert A\right)  \left(  A\right\vert \right\} \\
&  =\operatorname{Tr}\left\{  \left[  \widehat{\Lambda}_{k},\left\vert
A\right)  \left(  A\right\vert \right]  \widehat{\Phi\left(  \varphi
_{j}\right)  ^{\dagger}\Phi\left(  \varphi_{j}\right)  }\right\} \\
&  =\operatorname{Tr}\left\{  \left[  \left\vert A\right)  \left(
A\right\vert ,\widehat{\Phi\left(  \varphi_{j}\right)  ^{\dagger}\Phi\left(
\varphi_{j}\right)  }\right]  \widehat{\Lambda}_{k}\right\}  ,
\end{align}
by considering%
\begin{equation}
\operatorname{Tr}\left\{  \left[  A,B\right]  D\right\}  =\operatorname{Tr}%
\left\{  ABD-BAD\right\}  .
\end{equation}%
\begin{equation}
\operatorname{Tr}\left\{  \widehat{\Lambda}_{k}\left(  A\right)  \Phi\left(
\varphi_{j}\right)  ^{\dagger}\Phi\left(  \varphi_{j}\right)  A-A\Phi\left(
\varphi_{j}\right)  ^{\dagger}\Phi\left(  \varphi_{j}\right)  \widehat
{\Lambda}_{k}\left(  A\right)  \right\}
\end{equation}
which is equivalent to
\begin{equation}
\delta\left(  A\left\vert \widehat{\Lambda}_{k}A\right.  \right)
=\delta\operatorname{Tr}\left\{  \left(  \widehat{\Lambda}_{k}\left(
A\right)  \right)  \Phi\left(  \varphi_{j}\right)  ^{\dagger}\Phi\left(
\varphi_{j}\right)  \left(  \widehat{\Lambda}_{k}\left(  A\right)  \right)
\right\}  =0
\end{equation}
wrt variations produced by the commutative subgroup, $\mathfrak{G}_{\varphi},$
of $U\left(  \mathcal{B}_{2}\left(  \mathcal{H}^{N}\right)  \right)  $ formed
by the unitary transformations
\begin{equation}
\exp i\left\{  \sum_{1\leq j\leq m}\lambda_{j}\widehat{\Phi\left(  \varphi
_{j}\right)  ^{\dagger}\Phi\left(  \varphi_{j}\right)  }\right\}  ;\lambda
_{j}\in\mathbb{R},1\leq j\leq m.\,
\end{equation}
$\lceil$%
\begin{equation}
\left\langle \Psi\left\vert \left[  \Phi\left(  \varphi_{j}\right)  ^{\dagger
}\Phi\left(  \varphi_{j}\right)  ,H\right]  \Psi\right.  \right\rangle =0
\end{equation}
The energy is stable wrt to one particle variations of the coefficients of the
state $\Psi$ in the basis $\left\{  \varphi_{j}\right\}  $ or equivalently the
one particle coefficients are stable to time variations.

Consider the CI expansion
\begin{equation}
\sum_{1\leq j_{1}<\cdots<j_{N}\leq r}C_{j_{1}\cdots j_{N}}\left\vert
\varphi_{j_{1}}\cdots\varphi_{j_{N}}\right\rangle
\end{equation}
then every state can be obtained by the general linear transformations
\begin{equation}
\exp\left\{  \sum_{1\leq M\leq N}\sum_{1\leq j_{1}<\cdots<j_{M}\leq r}%
Z_{j_{1}\cdots j_{M}}\Phi\left(  \varphi_{j_{1}}\right)  ^{\dagger}\Phi\left(
\varphi_{j_{1}}\right)  \cdots\Phi\left(  \varphi_{j_{M}}\right)  ^{\dagger
}\Phi\left(  \varphi_{j_{M}}\right)  \right\}
\end{equation}
$\rfloor$

The set of operators $\left\{  \left.  \widehat{\Lambda}_{j}\right\vert 1\leq
j\leq\nu\right\}  $ produces a basis $\left\{  \left.  \left\vert
\widehat{\Lambda}_{j}A\right)  \right\vert 1\leq j\leq\nu\right\}  $ of the
tangent space$.\mathfrak{G}_{\varphi}$ is a subgroup of $\mathfrak{G}.$

Clearly a subgroup $\mathfrak{G}$ of transformations that satisfy this
condition are produced by
\begin{align}
U  &  =e^{i\Lambda}\\
\Lambda &  \in\operatorname{alg}\left\{  \left.  \Phi\left(  \varphi
_{j}\right)  ^{\dagger}\Phi\left(  \varphi_{j}\right)  \right\vert 1\leq j\leq
m\right\} \\
\widehat{U}\left(  A\right)   &  =U^{\dagger}AU,
\end{align}
another subgroup is given by the inner isometry group of $A$ i.e.%
\begin{equation}
U^{\dagger}AU=A.
\end{equation}


\subsection{General Function}

The kernel must be at least $n-m$ dimensional as
\begin{equation}
d\boldsymbol{g}\left(  \boldsymbol{A}\right)  :\mathbb{R}^{n}\rightarrow
\mathbb{R}^{m}%
\end{equation}
and the equation
\[
d\boldsymbol{g}\left(  \boldsymbol{A}\right)  \left(  \delta\boldsymbol{A}%
\right)  =\boldsymbol{0}%
\]
must have at least $n-m$ linearly independent solutions. It might have more
solutions, which could produce a larger kernel.

One can decompose $\mathbb{R}^{n}$ into the direct sum
\[
\mathbb{R}^{n}=\ker\left\{  d\boldsymbol{g}\left(  \boldsymbol{A}\right)
\right\}  \oplus\ker\left\{  d\boldsymbol{g}\left(  \boldsymbol{A}\right)
\right\}  ^{\bot}%
\]
and
\[
\boldsymbol{dg}\left(  \boldsymbol{A}\right)  :\ker\left\{  d\boldsymbol{g}%
\left(  \boldsymbol{A}\right)  \right\}  ^{\bot}\overset{1-1}{\rightarrow
}\operatorname*{Range}\left\{  d\boldsymbol{g}\left(  \boldsymbol{A}\right)
\right\}  ,
\]
i.e. $\ker\left\{  d\boldsymbol{g}\left(  \boldsymbol{A}\right)  \right\}
^{\bot}\cong\operatorname*{Range}\left\{  d\boldsymbol{g}\left(
\boldsymbol{A}\right)  \right\}  $.

$\lceil$%
\begin{align}
\mathbf{x}  &  \mathbf{\in}\ker\left\{  d\boldsymbol{g}\left(  \boldsymbol{A}%
\right)  \right\}  \equiv\left(  \mathbf{y}\left\vert d\boldsymbol{g}\left(
\boldsymbol{A}\right)  \mathbf{x}\right.  \right)  =0\forall
\mathbf{y\Rightarrow}\left(  d\boldsymbol{g}\left(  \boldsymbol{A}\right)
^{\dagger}\mathbf{y}\left\vert \mathbf{x}\right.  \right)  =0\nonumber\\
&  \Rightarrow d\boldsymbol{g}\left(  \boldsymbol{A}\right)  ^{\dagger
}\mathbf{y\bot}\ker\left\{  d\boldsymbol{g}\left(  \boldsymbol{A}\right)
\right\}  \text{ i.e. }\operatorname*{Range}\left\{  d\boldsymbol{g}\left(
\boldsymbol{A}\right)  ^{\dagger}\right\}  \mathbf{\bot}\ker\left\{
d\boldsymbol{g}\left(  \boldsymbol{A}\right)  \right\}
\end{align}
$\rfloor$

The \emph{normal space, }$\mathcal{N}_{\boldsymbol{A}}\boldsymbol{g},$ of
$\boldsymbol{g}$ at the point $\boldsymbol{A}$ is%
\[
\mathcal{N}_{\boldsymbol{A}}\boldsymbol{g}=\left(  \mathcal{T}_{\boldsymbol{A}%
}\right)  ^{\bot},
\]
and
\begin{equation}
\operatorname*{Range}\left\{  d\boldsymbol{g}\left(  \boldsymbol{A}\right)
\right\}  \cong\ker\left\{  d\boldsymbol{g}\left(  \boldsymbol{A}\right)
\right\}  ^{\bot}\subseteq\mathcal{N}_{\boldsymbol{A}}\boldsymbol{g}%
\subseteq\mathbb{R}^{m}.
\end{equation}
Thus there might be normals outside the range of $d\boldsymbol{g}\left(
\boldsymbol{A}\right)  $.

\subsection{General Dual Spaces}

Observe that
\[
\boldsymbol{g}:V\longrightarrow W\equiv\boldsymbol{g}:\mathbb{R}%
^{n}\longrightarrow\mathbb{R}^{m}%
\]
$\lceil$In the DFT example
\begin{align}
V  &  =V^{d}=\mathcal{B}_{2sa}\left(  \mathcal{H}^{N}\right) \\
W  &  =L_{1}\left(  \Delta_{A}\right) \\
W^{\boldsymbol{d}}  &  =L_{\infty}\left(  \Delta_{A}\right) \\
\boldsymbol{g}\left(  A\right)   &  =\operatorname{Tr}\left\{  \Phi^{\dagger
}\Phi_{\delta}A^{2}\right\}
\end{align}
$\rfloor$

as theses are linear spaces we have the following relationships%
\[
d\boldsymbol{g}\left(  \boldsymbol{A}\right)  :V\longrightarrow W\equiv
d\boldsymbol{g}\left(  \boldsymbol{A}\right)  :\mathbb{R}^{n}\longrightarrow
\mathbb{R}^{m}%
\]
$\left[  \left[  d\boldsymbol{g}\left(  A\right)  \right]  \left(  \delta
A\right)  =2\operatorname{Tr}\left\{  A\Phi^{\dagger}\Phi_{\delta}\delta
A\right\}  \right]  $%
\begin{align*}
\left\langle \bullet,\bullet\right\rangle _{V} &  :V^{\boldsymbol{d}}\times
V\longrightarrow\mathbb{R}\\
\left\langle \bullet,\bullet\right\rangle _{W} &  :W^{\boldsymbol{d}}\times
W\longrightarrow\mathbb{R}%
\end{align*}
$\left[
\begin{array}
[c]{c}%
\begin{array}
[c]{c}%
\begin{array}
[c]{c}%
\left\langle \left.  \lambda\right\vert 2\operatorname{Tr}\left\{
A\Phi^{\dagger}\Phi_{\delta}\delta A\right\}  \right\rangle _{L_{\infty
}\left(  \Delta_{A}\right)  \times L_{1}\left(  \Delta_{A}\right)  }=\left(
\left.  \left[  d\boldsymbol{g}\left(  A\right)  \right]  ^{\dagger}\left(
\lambda\right)  \right\vert \delta A\right)  _{\mathcal{B}_{2sa}\left(
\mathcal{H}^{N}\right)  }\\
=2\int_{\Delta_{A}}\lambda\left(  \boldsymbol{r}\right)  \operatorname{Tr}%
\left\{  A\Phi^{\dagger}\left(  \boldsymbol{r}\right)  \Phi\left(
\boldsymbol{r}\right)  \delta A\right\}  d\boldsymbol{r}\\
=2\int_{\Delta_{A}}\operatorname{Tr}\left\{  A\lambda\left(  \boldsymbol{r}%
\right)  \Phi^{\dagger}\left(  \boldsymbol{r}\right)  \Phi\left(
\boldsymbol{r}\right)  dx\delta A\right\}  d\boldsymbol{r}=2\operatorname{Tr}%
\left\{  \mathfrak{l}_{A}\delta A\right\}  \\
=2\left(  \left.  \left[  A\mathbf{\lambda}\right]  _{+}\right\vert \delta
A\right)  _{\mathcal{B}_{2sa}\left(  \mathcal{H}^{N}\right)  }=2\left(
\left.  \mathfrak{l}_{A}\right\vert \delta A\right)  _{\mathcal{B}%
_{2sa}\left(  \mathcal{H}^{N}\right)  }%
\end{array}
\\
\mathfrak{l}_{A}=\left[  A\mathbf{\lambda}\right]  _{+}=\int_{\Delta_{A}%
}\lambda\left(  \boldsymbol{r}\right)  \left[  A\Phi^{\dagger}\left(
\boldsymbol{r}\right)  \Phi\left(  \boldsymbol{r}\right)  \right]
_{+}d\boldsymbol{r}=\frac{1}{2}\int_{\Delta_{A}}\lambda\left(  \boldsymbol{r}%
\right)  \left\{  A\Phi^{\dagger}\left(  \boldsymbol{r}\right)  \Phi\left(
\boldsymbol{r}\right)  +\Phi^{\dagger}\left(  \boldsymbol{r}\right)
\Phi\left(  \boldsymbol{r}\right)  A\right\}  d\boldsymbol{r}\\
\in\mathcal{B}_{2sa}\left(  \mathcal{H}^{N}\right)  ,\,\lambda\in L_{\infty
}\left(  \Delta_{A}\right)  \\
\mathbf{\lambda=}\int_{\Delta_{A}}\lambda\left(  \boldsymbol{r}\right)
\Phi^{\dagger}\left(  \boldsymbol{r}\right)  \Phi\left(  \boldsymbol{r}%
\right)  d\boldsymbol{r}%
\end{array}
\\
\mathcal{B}\left(  \mathcal{H}^{N}\right)  \supseteq\mathcal{C}\left(
\mathcal{H}^{N}\right)  \supseteq\mathcal{B}_{2sa}\left(  \mathcal{H}%
^{N}\right)  \supseteq\mathcal{B}_{1sa}\left(  \mathcal{H}^{N}\right)
\supseteq\mathcal{B}_{0sa}\left(  \mathcal{H}^{N}\right)
\end{array}
\right]  $

$\left[
\begin{array}
[c]{c}%
\int_{\Delta_{A}}\lambda\left(  \boldsymbol{r}\right)  \left\vert \left[
A\Phi^{\dagger}\left(  \boldsymbol{r}\right)  \Phi\left(  \boldsymbol{r}%
\right)  \right]  _{+}\right)  d\boldsymbol{r}=0\Longleftrightarrow
\lambda\left(  \boldsymbol{r}\right)  \overset{a.e.}{=}0\\
\int_{\Delta_{A}}\lambda\left(  \boldsymbol{r}\right)  \left(  \left[
A\Phi^{\dagger}\left(  \boldsymbol{r}\right)  \Phi\left(  \boldsymbol{r}%
\right)  \right]  _{+}\right\vert d\boldsymbol{r}=0\Longleftrightarrow
\lambda\left(  \boldsymbol{r}\right)  \overset{a.e.}{=}0\\
\ker\left\{  dg\left(  A\right)  ^{\dagger}\right\}  =\left\{  \lambda
\left\vert \left[  dg\left(  A\right)  ^{\dagger}\right]  \left(
\lambda\right)  =0\right.  \right\}  \\
\left[  dg\left(  A\right)  ^{\dagger}\right]  \left(  \lambda\right)
=0\Longleftrightarrow\left(  \left.  \left[  dg\left(  A\right)  ^{\dagger
}\right]  \left(  \lambda\right)  \right\vert \delta A\right)  =0\forall\delta
A\in\mathcal{B}_{2sa}\left(  \mathcal{H}^{N}\right)  \\
=\left\langle \left.  \lambda\right\vert \left[  dg\left(  A\right)  \right]
\left(  \delta A\right)  \right\rangle _{L_{\infty}\left(  \Delta_{A}\right)
\times L_{1}\left(  \Delta_{A}\right)  }=0\forall\delta A\in\mathcal{B}%
_{2sa}\left(  \mathcal{H}^{N}\right)  \\
\Rightarrow\lambda\bot\operatorname{range}\left\{  dg\left(  A\right)
\right\}  \equiv\lambda\in\left[  \operatorname{range}\left\{  dg\left(
A\right)  \right\}  \right]  ^{\bot}\text{ }\\
\text{i.e. }\lambda\in\ker\left\{  dg\left(  A\right)  ^{\dagger}\right\}
\Longleftrightarrow\lambda\in\left[  \operatorname{range}\left\{  dg\left(
A\right)  \right\}  \right]  ^{\bot}\\
\operatorname{range}\left\{  dg\left(  A\right)  \right\}  =L_{1}\left(
\Delta_{A}\right)  \text{ }\equiv dg\left(  A\right)  \text{ is surjective}\\
dg\left(  A\right)  \text{ is surjective }\equiv\left[  \operatorname{range}%
\left\{  dg\left(  A\right)  \right\}  \right]  ^{\bot}=\phi\\
\left(  \left.  \left[  dg\left(  A\right)  ^{\dagger}\right]  \left(
\lambda\right)  \right\vert \delta A\right)  =\left\langle \left.
\lambda\right\vert \left[  dg\left(  A\right)  \right]  \left(  \delta
A\right)  \right\rangle _{L_{\infty}\left(  \Delta_{A}\right)  \times
L_{1}\left(  \Delta_{A}\right)  }=0\forall\delta A\in\mathcal{B}_{2sa}\left(
\mathcal{H}^{N}\right)  \\
\Longleftrightarrow\int_{\Delta_{A}}\lambda\left(  \boldsymbol{r}\right)
\left(  \delta A\left\vert \left[  A\Phi^{\dagger}\left(  \boldsymbol{r}%
\right)  \Phi\left(  \boldsymbol{r}\right)  \right]  _{+}\right.  \right)
d\boldsymbol{r}=0\forall\delta A\in\mathcal{B}_{2sa}\left(  \mathcal{H}%
^{N}\right)  \\
\Longleftrightarrow\int_{\Delta_{A}}\lambda\left(  \boldsymbol{r}\right)
\left\vert \left[  A\Phi^{\dagger}\left(  \boldsymbol{r}\right)  \Phi\left(
\boldsymbol{r}\right)  \right]  _{+}\right)  d\boldsymbol{r}=0\\
\left[  \operatorname{range}\left\{  dg\left(  A\right)  \right\}  \right]
^{\bot}=\phi\equiv\int_{\Delta_{A}}\lambda\left(  \boldsymbol{r}\right)
\left\vert \left[  A\Phi^{\dagger}\left(  \boldsymbol{r}\right)  \Phi\left(
\boldsymbol{r}\right)  \right]  _{+}\right)  d\boldsymbol{r}%
=0\Longleftrightarrow\lambda\left(  \boldsymbol{r}\right)  \overset{a.e.}{=}0
\end{array}
\right]  $%
\begin{align*}
d\boldsymbol{g}\left(  \boldsymbol{A}\right)  ^{\dagger} &  :W^{\boldsymbol{d}%
}\longrightarrow V^{\boldsymbol{d}}\\
\left\langle \left.  d\boldsymbol{g}\left(  \boldsymbol{A}\right)  ^{\dagger
}\boldsymbol{w}_{\boldsymbol{d}}\right\vert \boldsymbol{v}\right\rangle _{V}
&  =\left\langle \left.  \boldsymbol{w}_{\boldsymbol{d}}\right\vert
d\boldsymbol{g}\left(  \boldsymbol{A}\right)  \boldsymbol{v}\right\rangle
_{W}\,\forall\boldsymbol{v\in}\operatorname*{domain}\left\{  d\boldsymbol{g}%
\left(  \boldsymbol{A}\right)  \right\}  \subseteq V
\end{align*}%
\[
\dagger:\mathcal{L}\left(  V,W\right)  \longrightarrow\mathcal{L}\left(
W^{\boldsymbol{d}},V^{\boldsymbol{d}}\right)
\]
$\lceil$%
\[
_{V}\left\langle d\boldsymbol{g}\left(  \boldsymbol{A}\right)  ^{\dagger
}\boldsymbol{w}_{\boldsymbol{A}}\right\vert =\,_{W}\left\langle \boldsymbol{w}%
_{\boldsymbol{A}}\right\vert d\boldsymbol{g}\left(  \boldsymbol{A}\right)
=\left(  d\boldsymbol{g}\left(  \boldsymbol{A}\right)  \left\vert
\boldsymbol{w}\right\rangle _{W}\right)  ^{\dagger}%
\]
if $W$ can be embedded in $W^{\boldsymbol{d}}\rfloor$

Hence%
\begin{align*}
\operatorname*{Range}\left\{  d\boldsymbol{g}\left(  \boldsymbol{A}\right)
\right\}  ^{\bot} &  =\left\{  \boldsymbol{w}_{\boldsymbol{d}}\in
W^{\boldsymbol{d}}\left\vert \left\langle \left.  \boldsymbol{w}%
_{\boldsymbol{d}}\right\vert \boldsymbol{w}\right\rangle _{W}=0\forall
\boldsymbol{w}\in\operatorname*{Range}\left\{  d\boldsymbol{g}\left(
\boldsymbol{A}\right)  \right\}  \subseteq W\right.  \right\}  \\
&  \equiv\left\{  \boldsymbol{w}_{\boldsymbol{d}}\in W^{\boldsymbol{d}%
}\left\vert \left\langle \left.  \boldsymbol{w}_{\boldsymbol{d}}\right\vert
\boldsymbol{w}\right\rangle _{W}=0\forall\boldsymbol{w}\in
\operatorname*{Range}\left\{  d\boldsymbol{g}\left(  \boldsymbol{A}\right)
\right\}  \subseteq\mathbb{R}^{m}\right.  \right\}
\end{align*}
and
\begin{align*}
\ker\left\{  d\boldsymbol{g}\left(  \boldsymbol{A}\right)  \right\}  ^{\bot}
&  =\left\{  \boldsymbol{v}_{\boldsymbol{A}}\in V^{\boldsymbol{d}}\left\vert
\left\langle \left.  \boldsymbol{v}_{\boldsymbol{d}}\right\vert \boldsymbol{v}%
\right\rangle _{V}=0\forall\boldsymbol{v}\in\ker\left\{  d\boldsymbol{g}%
\left(  \boldsymbol{A}\right)  \right\}  \subseteq V\right.  \right\}  \\
&  \equiv\left\{  \boldsymbol{v}_{\boldsymbol{A}}\in V^{\boldsymbol{d}%
}\left\vert \left\langle \left.  \boldsymbol{v}_{\boldsymbol{d}}\right\vert
\boldsymbol{v}\right\rangle _{V}=0\forall\boldsymbol{v}\in\ker\left\{
d\boldsymbol{g}\left(  \boldsymbol{A}\right)  \right\}  \subseteq
\mathbb{R}^{n}\right.  \right\}  .
\end{align*}


If $\operatorname*{Range}\left\{  d\boldsymbol{g}\left(  \boldsymbol{A}%
\right)  \right\}  =\mathbb{R}^{m}$ then $\left\{  \left.  \nabla g_{j}\left(
\boldsymbol{A}\right)  \right\vert 1\leq j\leq m\right\}  $ are linearly
independent and the vector $\left\langle \boldsymbol{v}\right\vert
\in\mathbb{R}^{m}$ can be uniquely expanded as
\begin{align}
\left\langle \boldsymbol{v}\right\vert  &  =\sum_{1\leq j\leq m}\lambda
_{j}\nabla g_{j}\left(  \boldsymbol{A}\right)  =\left(
\begin{array}
[c]{ccc}%
\lambda_{1} & ,\cdots, & \lambda_{m}%
\end{array}
\right)  \left(
\begin{array}
[c]{c}%
\nabla g_{1}\left(  \boldsymbol{A}\right) \\
\vdots\\
\nabla g_{m}\left(  \boldsymbol{A}\right)
\end{array}
\right) \nonumber\\
&  =\left(
\begin{array}
[c]{ccc}%
\lambda_{1} & ,\cdots, & \lambda_{m}%
\end{array}
\right)  \left(
\begin{array}
[c]{ccc}%
\frac{\partial g_{1}\left(  \boldsymbol{A}\right)  }{\partial\boldsymbol{A}%
_{1}} & \cdots & \frac{\partial g_{1}\left(  \boldsymbol{A}\right)  }%
{\partial\boldsymbol{A}_{n}}\\
\vdots & \ddots & \vdots\\
\frac{\partial g_{m}\left(  \boldsymbol{A}\right)  }{\partial\boldsymbol{A}%
_{1}} & \cdots & \frac{\partial g_{m}\left(  \boldsymbol{A}\right)  }%
{\partial\boldsymbol{A}_{n}}%
\end{array}
\right) \nonumber\\
&  =\left\langle \boldsymbol{\lambda}\right\vert \nabla\boldsymbol{g}\left(
\boldsymbol{A}\right)  \equiv\left\langle \boldsymbol{\lambda}\right\vert
d\boldsymbol{g}\left(  \boldsymbol{A}\right)  .
\end{align}
If $\operatorname*{Range}\left\{  d\boldsymbol{g}\left(  \boldsymbol{A}%
\right)  \right\}  \subset\mathbb{R}^{m}$ then $d\boldsymbol{g}\left(
\boldsymbol{A}\right)  $ has a non empty \emph{left kernel} i.e.
\begin{equation}
\exists\,\left\langle \boldsymbol{v}\right\vert \backepsilon\left\langle
\boldsymbol{\lambda}\right\vert d\boldsymbol{g}\left(  \boldsymbol{A}\right)
=\boldsymbol{0}\Longleftrightarrow\boldsymbol{\lambda\in}\ker\left\{
d\boldsymbol{g}\left(  \boldsymbol{A}\right)  ^{\dagger}\right\}
\end{equation}
and%
\begin{align}
\operatorname*{Range}\left\{  d\boldsymbol{g}\left(  \boldsymbol{A}\right)
\right\}  ^{\bot}  &  =\left\{  \left.  \boldsymbol{\lambda}\right\vert
\left\langle \left.  \boldsymbol{\lambda}\right\vert \left[  d\boldsymbol{g}%
\left(  \boldsymbol{A}\right)  \right]  \left(  \delta\boldsymbol{A}\right)
\right\rangle =0\forall\delta\boldsymbol{A}\in\operatorname*{domain}\left\{
d\boldsymbol{g}\left(  \boldsymbol{A}\right)  \right\}  \right\} \nonumber\\
&  =\left\{  \left.  \boldsymbol{\lambda}\right\vert \left\langle \left.
\boldsymbol{\lambda}\right\vert \left[  d\boldsymbol{g}\left(  \boldsymbol{A}%
\right)  \right]  \left(  \delta\boldsymbol{A}\right)  \right\rangle
=0\forall\delta\boldsymbol{A}\in\ker\left\{  d\boldsymbol{g}\left(
\boldsymbol{A}\right)  \right\}  ^{\bot}\right\}
\end{align}
If $\boldsymbol{\eta}\neq\boldsymbol{0\in}\ker\left\{  d\boldsymbol{g}\left(
\boldsymbol{A}\right)  ^{\dagger}\right\}  $ then
\begin{align}
\left\langle \left.  \boldsymbol{\eta}\right\vert \left[  d\boldsymbol{g}%
\left(  \boldsymbol{A}\right)  \right]  \left(  \delta\boldsymbol{A}\right)
\right\rangle  &  =0\forall\delta\boldsymbol{A}\in V=\left\langle \left.
\boldsymbol{\eta}\right\vert \boldsymbol{w}\right\rangle =0\forall
\boldsymbol{w\in}\operatorname*{Range}\left\{  d\boldsymbol{g}\left(
\boldsymbol{A}\right)  \right\} \nonumber\\
&  \Rightarrow\operatorname*{Range}\left\{  d\boldsymbol{g}\left(
\boldsymbol{A}\right)  \right\}  \subset W
\end{align}
thus $d\boldsymbol{g}\left(  \boldsymbol{A}\right)  $ is not surjective.
Further
\begin{equation}
\left\langle \boldsymbol{\lambda}\right\vert \left[  d\boldsymbol{g}\left(
\boldsymbol{A}\right)  \right]  =\left\langle \boldsymbol{\lambda
}+\boldsymbol{\eta}\right\vert \left[  d\boldsymbol{g}\left(  \boldsymbol{A}%
\right)  \right]  \text{ \thinspace for }\boldsymbol{\eta}\in\ker\left\{
d\boldsymbol{g}\left(  \boldsymbol{A}\right)  ^{\dagger}\right\}
\end{equation}
so functionals on $\operatorname*{Range}\left\{  d\boldsymbol{g}\left(
\boldsymbol{A}\right)  \right\}  $ are not unique.

\subsection{Density Functional Example}

For the case $\boldsymbol{f}\left(  A\right)  $ one has in the infinite
dimensional case $V=V^{d}=\mathcal{B}_{sa}\left(  \mathcal{H}^{N}\right)  ,$
$W=L_{1}\left(  \mathbb{R}^{3}\right)  $ and $W^{d}=L_{\infty}\left(
\mathbb{R}^{3}\right)  ,$while in the finite dimensional case $V=V^{d}%
=\mathbb{R}^{n},$ and $V=V^{d}=\mathbb{R}^{m}.$

\subsection{General Function}%

\begin{align}
\mathbb{R}^{n}  &  =\ker\left\{  d\boldsymbol{g}\left(  \boldsymbol{A}\right)
\right\}  ^{\bot}\oplus\ker\left\{  d\boldsymbol{g}\left(  \boldsymbol{A}%
\right)  \right\} \nonumber\\
\mathbb{R}^{m}  &  =\operatorname*{Range}\left\{  d\boldsymbol{g}\left(
\boldsymbol{A}\right)  \right\}  ^{\bot}\oplus\operatorname*{Range}\left\{
d\boldsymbol{g}\left(  \boldsymbol{A}\right)  \right\} \nonumber\\
\operatorname*{Range}\left\{  d\boldsymbol{g}\left(  \boldsymbol{A}\right)
\right\}   &  \cong\ker\left\{  d\boldsymbol{g}\left(  \boldsymbol{A}\right)
\right\}  ^{\bot}%
\end{align}
and in all cases
\begin{equation}
\mathcal{T}_{\boldsymbol{A}}\subseteq\ker\left\{  d\boldsymbol{g}\left(
\boldsymbol{A}\right)  \right\}  .
\end{equation}
A point on the constraint manifold is called \emph{normal }if $\left\{
\left.  \nabla g_{j}\left(  \boldsymbol{A}\right)  \right\vert 1\leq j\leq
m\right\}  $ generates the normal space (thus is a linearly independent basis
or generalized basis for this space)
\begin{equation}
\text{in general }\operatorname*{Range}\left\{  d\boldsymbol{g}\left(
\boldsymbol{A}\right)  \right\}  \subseteq\mathcal{N}_{\boldsymbol{A}%
}\boldsymbol{g}\Rightarrow\operatorname*{Range}\left\{  d\boldsymbol{g}\left(
\boldsymbol{A}\right)  \right\}  ^{\bot}\supseteq\left(  \mathcal{N}%
_{\boldsymbol{A}}\boldsymbol{g}\right)  ^{\bot}%
\end{equation}
and \emph{regular }if $\mathcal{T}_{\boldsymbol{A}}=\ker\left\{
d\boldsymbol{g}\left(  \boldsymbol{A}\right)  \right\}  $%
\begin{equation}
\text{in general }\mathcal{T}_{\boldsymbol{A}}\subseteq\ker\left\{
d\boldsymbol{g}\left(  \boldsymbol{A}\right)  \right\}  \Rightarrow\left(
\mathcal{T}_{\boldsymbol{A}}\right)  ^{\bot}\supseteq\ker\left\{
d\boldsymbol{g}\left(  \boldsymbol{A}\right)  \right\}  ^{\bot}.
\end{equation}
\emph{Normal implies regular. }Thus for a normal point $\boldsymbol{A}$
\begin{align}
\mathcal{N}_{\boldsymbol{A}}\boldsymbol{g}  &  =\left(  \mathcal{T}%
_{\boldsymbol{A}}\right)  ^{\bot}=\ker\left\{  d\boldsymbol{g}\left(
\boldsymbol{A}\right)  \right\}  ^{\bot}\nonumber\\
\ker\left\{  d\boldsymbol{g}\left(  \boldsymbol{A}\right)  \right\}  ^{\bot}
&  \cong\operatorname*{Range}\left\{  d\boldsymbol{g}\left(  \boldsymbol{A}%
\right)  \right\}  =W.
\end{align}


e.g. 1%

\begin{align*}
&  g:\mathbb{R}^{2}\rightarrow\mathbb{R}\\
&  g\left(  x,y\right)  =x^{2}+y^{2}=c^{2};\quad y=f\left(  x\right)
=+\sqrt{c^{2}-x^{2}};\quad\\
&  \left[  df\left(  x\right)  \right]  \left(  \delta x\right)  =\left[
-x\left(  c^{2}-x^{2}\right)  ^{-\frac{1}{2}}\right]  \left(  \delta x\right)
\\
&  \left\langle \nabla g\left(  x,y\right)  \right\vert =\left(  2x,2y\right)
\\
&  \ker\left\{  dg\left(  x,y\right)  \right\}  =\left\{  \left.
\alpha\left(
\begin{array}
[c]{c}%
y\\
-x
\end{array}
\right)  \right\vert \alpha\in\mathbb{R}\right\}  \cong\mathbb{R}\\
&  x\left(  t\right)  =c\cos\left(  \omega t+\omega_{0}\right)  ,\,y\left(
t\right)  =c\sin\left(  \omega t+\omega_{0}\right) \\
\theta\left(  t\right)   &  =\omega t+\omega_{0}=\tan^{-1}\left\{
\frac{y\left(  t\right)  }{x\left(  t\right)  }\right\}  ;\omega\in
\mathbb{R},\quad\tan^{-1}\text{ is the principal value }\\
&  \text{for }x\neq0\text{ then }\\
\theta\left(  0\right)   &  =\omega_{0}=\left\{
\begin{array}
[c]{c}%
\tan^{-1}\left\{  \frac{y}{x}\right\}  ;x>0\\
\tan^{-1}\left\{  \frac{y}{x}\right\}  -\pi;x<0
\end{array}
\right. \\
\text{for }x  &  =0\text{ then }\omega_{0}=\left\{
\begin{array}
[c]{c}%
\frac{\pi}{2};y>0\\
-\frac{\pi}{2};y<0
\end{array}
\right. \\
&  \mathcal{T}_{\left(  x,y\right)  }g=\left(
\begin{array}
[c]{c}%
\delta x\\
\delta y
\end{array}
\right)  =\left\{  \omega\left(
\begin{array}
[c]{c}%
-c\sin\left(  \omega_{0}\right) \\
c\cos\left(  \omega_{0}\right)
\end{array}
\right)  =\omega\left(
\begin{array}
[c]{c}%
-y\\
x
\end{array}
\right)  \right. \\
&  \left.  ,\omega\in\mathbb{R}\equiv\alpha\left(
\begin{array}
[c]{c}%
y\\
-x
\end{array}
\right)  \right\} \\
&  \mathcal{N}_{\left(  x,y\right)  }g=\left[  \mathcal{T}_{\left(
x,y\right)  }g\right]  ^{\bot}\cong\mathbb{R}%
\end{align*}


e.g. 2%

\begin{align*}
g  &  :\mathbb{R}^{3}\rightarrow\mathbb{R};\quad g\left(  x,y,z\right)
=x^{2}+y^{2}+z^{2}-c^{2}=0;\quad dg\left(  x,y,z\right)  :\mathbb{R}%
^{3}\rightarrow\mathbb{R}\\
&  \left[  dg\left(  x,y,z\right)  \right]  \left(  \delta x,\delta y,\delta
z\right)  =\left(  2x,2y,2z\right)  \left(
\begin{array}
[c]{c}%
\delta x\\
\delta y\\
\delta z
\end{array}
\right)  \equiv\left\langle \left.  \nabla g\left(  x,y,z\right)  \right\vert
\delta\boldsymbol{x}\right\rangle \\
z  &  =f\left(  x,y\right)  =+\sqrt{c^{2}-x^{2}-y^{2}}\quad f:\mathbb{R}%
^{2}\rightarrow\mathbb{R};\quad df\left(  x,y\right)  :\mathbb{R}%
^{2}\rightarrow\mathbb{R}\\
&  \left[  df\left(  x,y\right)  \right]  \left(  \delta x,\delta y\right)
=\frac{-1}{\sqrt{c^{2}-x^{2}-y^{2}}}\left(  x,y\right)  \left(
\begin{array}
[c]{c}%
\delta x\\
\delta y
\end{array}
\right)  \equiv\left\langle \left.  \nabla f\left(  x,y\right)  \right\vert
\delta\boldsymbol{x}\right\rangle \\
&  g\left(  x,y,z\right)  =F\left(  \left(  x,y\right)  ,f\left(  x,y\right)
\right)  =0=x^{2}+y^{2}+z^{2}-c^{2}\equiv\operatorname*{graph}\left\{  \text{
}f\right\} \\
&  \ker\left\{  dg\left(  x,y,z\right)  \right\}  =\left\{  \left.  \left(
\begin{array}
[c]{c}%
\delta x\\
\delta y\\
\delta z
\end{array}
\right)  \right\vert \left(  2x,2y,2z\right)  \left(
\begin{array}
[c]{c}%
\delta x\\
\delta y\\
\delta z
\end{array}
\right)  =0\right\}  \cong\mathbb{R}^{2}\\
&  \mathcal{N}_{\left(  x,y,z\right)  }g=\left\{  \left.  \alpha\left(
x,y,z\right)  \right\vert \alpha\in\mathbb{R}\right\}  \cong\mathbb{R}%
\end{align*}


e.g. 3%

\begin{align*}
&  g,dg\left(  \boldsymbol{r}\right)  :\mathbb{R}^{3}\rightarrow\mathbb{R}%
^{2}\equiv V\longrightarrow W;\,dg\left(  \boldsymbol{r}\right)  ^{\dagger
}:\mathbb{R}^{2}\rightarrow\mathbb{R}^{3}\equiv V^{\boldsymbol{A}%
}\longrightarrow W^{\boldsymbol{A}}\\
&  g_{1}\left(  x,y,z\right)  =x^{2}+y^{2}+z^{2}=c_{1}^{2}\\
&  g_{2}\left(  x,y,z\right)  =x^{2}+y^{2}=c_{2}^{2}\\
&  \left\langle \nabla g_{1}\left(  x,y,z\right)  \right\vert =\left(
2x,2y,2z\right) \\
&  \left\langle \nabla g_{2}\left(  x,y,z\right)  \right\vert =\left(
2x,2y,0\right) \\
&  \ker\left\{  d\boldsymbol{g}\left(  x,y,z\right)  \right\}  =\left\{
\left.  \left(
\begin{array}
[c]{c}%
\delta x\\
\delta y\\
\delta z
\end{array}
\right)  \right\vert \left(
\begin{array}
[c]{ccc}%
2x & 2y & 2z\\
2x & 2y & 0
\end{array}
\right)  \left(
\begin{array}
[c]{c}%
\delta x\\
\delta y\\
\delta z
\end{array}
\right)  =\left(
\begin{array}
[c]{c}%
0\\
0
\end{array}
\right)  \right\}  \subset V\equiv\mathbb{R}^{3}\\
&  \left.
\begin{array}
[c]{c}%
2x\delta x+2y\delta y+2z\delta z=0\\
2x\delta x+2y\delta y=0
\end{array}
\right\}  \Rightarrow\left.
\begin{array}
[c]{c}%
z\delta z=0\Rightarrow\delta z=0\text{ if }z\neq0\text{ and }\delta
z\in\mathbb{R}\text{ if }z=0\\
\delta x=-\frac{y}{x}\delta y\text{ or }\delta y=-\frac{x}{y}\delta x
\end{array}
\right\} \\
&  \ker\left\{  d\boldsymbol{g}\left(  x,y,z\right)  \right\}  =\left\{
\left.  \left(
\begin{array}
[c]{c}%
\delta x\\
-\frac{x}{y}\delta x\\
0
\end{array}
\right)  \right\vert \left(
\begin{array}
[c]{ccc}%
x & y & z
\end{array}
\right)  \neq\left(
\begin{array}
[c]{ccc}%
0 & 0 & 0
\end{array}
\right)  \right\}
\end{align*}
The kernel of the constraint
\[
\ker\left\{  d\boldsymbol{g}\left(  x,y,z\right)  \right\}  ==\left\{  \left.
\left(
\begin{array}
[c]{c}%
\delta x\\
\delta y\\
\delta z
\end{array}
\right)  \right\vert \left(
\begin{array}
[c]{ccc}%
2x & 2y & 2z\\
2x & 2y & 0
\end{array}
\right)  \left(
\begin{array}
[c]{c}%
\delta x\\
\delta y\\
\delta z
\end{array}
\right)  =\left(
\begin{array}
[c]{c}%
0\\
0
\end{array}
\right)  \right\}  \subset V\equiv\mathbb{R}^{3}%
\]%
\[
\left.
\begin{array}
[c]{c}%
2x\delta x+2y\delta y+2z\delta z=0\\
2x\delta x+2y\delta y=0
\end{array}
\right\}  \Rightarrow\left.
\begin{array}
[c]{c}%
z\delta z=0\Rightarrow\delta z=0\text{ if }z\neq0\text{ and }\delta
z\in\mathbb{R}\text{ if }z=0\\
\delta x=-\frac{y}{x}\delta y\text{ or }\delta y=-\frac{x}{y}\delta x
\end{array}
\right\}
\]%
\begin{align*}
x  &  =0,y\neq0\\
&
\begin{array}
[c]{c}%
2y\delta y+2z\delta z=0\\
2y\delta y=0\Rightarrow\delta y=0\Rightarrow2z\delta z=0
\end{array}
\end{align*}%
\begin{align*}
\ker\left\{  d\boldsymbol{g}\left(  x,y,z\right)  \right\}   &  =\left\{
\left.  \alpha\left(
\begin{array}
[c]{c}%
1\\
-\frac{x}{y}\\
0
\end{array}
\right)  \right\vert \alpha\in\mathbb{R}\right\} \\
\ker\left\{  d\boldsymbol{g}\left(  x,y,0\right)  \right\}   &  =\left\{
\left.  \alpha\left(
\begin{array}
[c]{c}%
1\\
-\frac{x}{y}\\
0
\end{array}
\right)  ,\beta\left(
\begin{array}
[c]{c}%
0\\
0\\
1
\end{array}
\right)  \right\vert \alpha,\beta\in\mathbb{R}\right\}
\end{align*}%
\begin{align*}
\ker\left\{  d\boldsymbol{g}\left(  0,y,z\right)  \right\}   &  =\left\{
\left.  \alpha\left(
\begin{array}
[c]{c}%
1\\
0\\
0
\end{array}
\right)  \right\vert \alpha\in\mathbb{R}\right\} \\
\ker\left\{  d\boldsymbol{g}\left(  0,y,0\right)  \right\}   &  =\left\{
\left.  \alpha\left(
\begin{array}
[c]{c}%
1\\
0\\
0
\end{array}
\right)  ,\beta\left(
\begin{array}
[c]{c}%
0\\
0\\
1
\end{array}
\right)  \right\vert \alpha,\beta\in\mathbb{R}\right\}
\end{align*}%
\begin{align*}
\ker\left\{  d\boldsymbol{g}\left(  x,0,z\right)  \right\}   &  =\left\{
\left.  \alpha\left(
\begin{array}
[c]{c}%
0\\
1\\
0
\end{array}
\right)  \right\vert \alpha\in\mathbb{R}\right\} \\
\ker\left\{  d\boldsymbol{g}\left(  x,0,0\right)  \right\}   &  =\left\{
\left.  \alpha\left(
\begin{array}
[c]{c}%
0\\
1\\
0
\end{array}
\right)  ,\beta\left(
\begin{array}
[c]{c}%
0\\
0\\
1
\end{array}
\right)  \right\vert \alpha,\beta\in\mathbb{R}\right\}
\end{align*}%
\begin{align*}
\ker\left\{  d\boldsymbol{g}\left(  0,0,z\right)  \right\}   &  =\left\{
\left.  \alpha\left(
\begin{array}
[c]{c}%
0\\
1\\
0
\end{array}
\right)  ,\beta\left(
\begin{array}
[c]{c}%
1\\
0\\
0
\end{array}
\right)  \right\vert \alpha,\beta\in\mathbb{R}\right\} \\
\ker\left\{  d\boldsymbol{g}\left(  0,0,0\right)  \right\}   &  =\left\{
\left.  \alpha\left(
\begin{array}
[c]{c}%
0\\
1\\
0
\end{array}
\right)  ,\beta\left(
\begin{array}
[c]{c}%
1\\
0\\
0
\end{array}
\right)  ,\gamma\left(
\begin{array}
[c]{c}%
0\\
0\\
1
\end{array}
\right)  \right\vert \alpha,\beta,\gamma\in\mathbb{R}\right\}
\end{align*}
The Normal space $\mathcal{N}_{\left(  x,y,z\right)  }\boldsymbol{g:}$ in
order to see if the bases are lin. indep we consider%
\begin{align*}
&  \alpha_{1}\left(
\begin{array}
[c]{ccc}%
2x & 2y & 2z
\end{array}
\right)  +\alpha_{2}\left(
\begin{array}
[c]{ccc}%
2x & 2y & 0
\end{array}
\right)  =0\\
&  \Rightarrow\alpha_{1}+\alpha_{2}=0\text{ if }x\neq0\text{ or }y\neq0\text{
and }\alpha_{1}z=0\\
&  \Rightarrow\alpha_{1}=\alpha_{2}=0\text{ if }z\neq0\text{ so vectors are
lin. indep. }\\
&  \Rightarrow\alpha_{1}=-\alpha_{2}\text{ if }z=0\text{ so vectors are lin.
dep.}\\
\mathcal{N}_{\left(  x,y,z\right)  }\boldsymbol{g}  &
=\operatorname*{linearspan}\left\{  \left(  x,y,z\right)  \right\}
=\operatorname*{range}\left\{  d\boldsymbol{g}\left(  \boldsymbol{r}\right)
\right\}  \subset W
\end{align*}%
\begin{align*}
\mathcal{N}_{\left(  x,y,z\right)  }\boldsymbol{g}  &
=\operatorname*{linearspan}\left\{  \left(  x,y,z\right)  ,\left(
x,y,0\right)  \right\}  =\operatorname*{range}\left\{  d\boldsymbol{g}\left(
\boldsymbol{r}\right)  \right\} \\
\mathcal{N}_{\left(  x,y,0\right)  }\boldsymbol{g}  &
=\operatorname*{linearspan}\left\{  \left(  x,y,0\right)  \right\}
=\operatorname*{range}\left\{  d\boldsymbol{g}\left(  \boldsymbol{r}\right)
\right\} \\
\mathcal{N}_{\left(  0,y,z\right)  }\boldsymbol{g}  &
=\operatorname*{linearspan}\left\{  \left(  0,y,z\right)  ,\left(
0,y,0\right)  \right\}  =\operatorname*{range}\left\{  d\boldsymbol{g}\left(
\boldsymbol{r}\right)  \right\} \\
\mathcal{N}_{\left(  0,y,0\right)  }\boldsymbol{g}  &
=\operatorname*{linearspan}\left\{  \left(  0,y,0\right)  \right\}
=\operatorname*{range}\left\{  d\boldsymbol{g}\left(  \boldsymbol{r}\right)
\right\} \\
\mathcal{N}_{\left(  x,0,z\right)  }\boldsymbol{g}  &
=\operatorname*{linearspan}\left\{  \left(  x,0,y\right)  ,\left(
x,0,0\right)  \right\}  =\operatorname*{range}\left\{  d\boldsymbol{g}\left(
\boldsymbol{r}\right)  \right\} \\
\mathcal{N}_{\left(  x,0,0\right)  }\boldsymbol{g}  &
=\operatorname*{linearspan}\left\{  \left(  x,0,0\right)  \right\}
=\operatorname*{range}\left\{  d\boldsymbol{g}\left(  \boldsymbol{r}\right)
\right\} \\
\mathcal{N}_{\left(  0,0,z\right)  }\boldsymbol{g}  &
=\operatorname*{linearspan}\left\{  \left(  0,0,z\right)  \right\}
=\operatorname*{range}\left\{  d\boldsymbol{g}\left(  \boldsymbol{r}\right)
\right\} \\
\mathcal{N}_{\left(  0,0,0\right)  }\boldsymbol{g}  &  =\left\{  \phi\right\}
=\operatorname*{range}\left\{  d\boldsymbol{g}\left(  \boldsymbol{r}\right)
\right\}
\end{align*}%
\[
V^{\boldsymbol{A}}\equiv\mathbb{R}^{3};W^{\boldsymbol{A}}\equiv\mathbb{R}^{2}%
\]%
\[
dg\left(  \boldsymbol{r}\right)  ^{\dagger}:W^{\boldsymbol{A}}\rightarrow
V^{\boldsymbol{A}}%
\]
Consider the Single Valued Decomposition (SVD)
\begin{equation}
\left(
\begin{array}
[c]{cc}%
u_{11} & u_{12}\\
u_{21} & u_{22}%
\end{array}
\right)  \left(
\begin{array}
[c]{ccc}%
2x & 2y & 2z\\
2x & 2y & 0
\end{array}
\right)  \left(
\begin{array}
[c]{ccc}%
v_{11} & v_{12} & v_{13}\\
v_{21} & v_{22} & v_{23}\\
v_{31} & v_{32} & v_{33}%
\end{array}
\right)  ^{\dagger}=\left(
\begin{array}
[c]{ccc}%
\alpha_{1} & 0 & 0\\
0 & \alpha_{2} & 0
\end{array}
\right)
\end{equation}%
\begin{equation}
\left(
\begin{array}
[c]{ccc}%
2x & 2y & 2z\\
2x & 2y & 0
\end{array}
\right)  =\left(
\begin{array}
[c]{cc}%
u_{11} & u_{12}\\
u_{21} & u_{22}%
\end{array}
\right)  ^{\dagger}\left(
\begin{array}
[c]{ccc}%
\alpha_{1} & 0 & 0\\
0 & \alpha_{2} & 0
\end{array}
\right)  \left(
\begin{array}
[c]{ccc}%
v_{11} & v_{12} & v_{13}\\
v_{21} & v_{22} & v_{23}\\
v_{31} & v_{32} & v_{33}%
\end{array}
\right)
\end{equation}%
\begin{align}
&  \left(
\begin{array}
[c]{ccc}%
2x & 2y & 2z\\
2x & 2y & 0
\end{array}
\right) \nonumber\\
&  =\alpha_{1}\left\vert \left(
\begin{array}
[c]{cc}%
w_{11} & w_{12}%
\end{array}
\right)  \right\rangle \left\langle \left(
\begin{array}
[c]{ccc}%
v_{11}^{\boldsymbol{A}} & v_{12}^{\boldsymbol{A}} & v_{13}^{\boldsymbol{A}}%
\end{array}
\right)  \right\vert +\alpha_{2}\left\vert \left(
\begin{array}
[c]{cc}%
w_{21} & w_{22}%
\end{array}
\right)  \right\rangle \left\langle \left(
\begin{array}
[c]{ccc}%
v_{21}^{\boldsymbol{A}} & v_{22}^{\boldsymbol{A}} & v_{23}^{\boldsymbol{A}}%
\end{array}
\right)  \right\vert
\end{align}%
\begin{align}
&  \left(
\begin{array}
[c]{ccc}%
2x & 2y & 2z\\
2x & 2y & 0
\end{array}
\right)  ^{\dagger}=\left(
\begin{array}
[c]{cc}%
2x & 2x\\
2y & 2y\\
2z & 0
\end{array}
\right)  \equiv dg\left(  \boldsymbol{r}\right)  ^{\dagger}\nonumber\\
&  =\alpha_{1}\left\vert \left(
\begin{array}
[c]{ccc}%
v_{11} & v_{12} & v_{13}%
\end{array}
\right)  \right\rangle \left\langle \left(
\begin{array}
[c]{cc}%
w_{11}^{\boldsymbol{A}} & w_{12}^{\boldsymbol{A}}%
\end{array}
\right)  \right\vert +\alpha_{2}\left\vert \left(
\begin{array}
[c]{ccc}%
v_{21} & v_{22} & v_{23}%
\end{array}
\right)  \right\rangle \left\langle \left(
\begin{array}
[c]{cc}%
w_{21}^{\boldsymbol{A}} & w_{22}^{\boldsymbol{A}}%
\end{array}
\right)  \right\vert
\end{align}%
\begin{align}
&  dg\left(  \boldsymbol{r}\right)  dg\left(  \boldsymbol{r}\right)
^{\dagger}=\left(
\begin{array}
[c]{ccc}%
2x & 2y & 2z\\
2x & 2y & 0
\end{array}
\right)  \left(
\begin{array}
[c]{ccc}%
2x & 2y & 2z\\
2x & 2y & 0
\end{array}
\right)  ^{\dagger}=\left(
\begin{array}
[c]{ccc}%
2x & 2y & 2z\\
2x & 2y & 0
\end{array}
\right)  \left(
\begin{array}
[c]{cc}%
2x & 2x\\
2y & 2y\\
2z & 0
\end{array}
\right) \nonumber\\
&  =\left(
\begin{array}
[c]{cc}%
4x^{2}+4y^{2}+4z^{2} & 4x^{2}+4y^{2}\\
4x^{2}+4y^{2} & 4x^{2}+4y^{2}%
\end{array}
\right) \nonumber\\
&  =\alpha_{1}^{2}\left\vert \left(
\begin{array}
[c]{cc}%
w_{11} & w_{12}%
\end{array}
\right)  \right\rangle \left\langle \left(
\begin{array}
[c]{cc}%
w_{11}^{\boldsymbol{A}} & w_{12}^{\boldsymbol{A}}%
\end{array}
\right)  \right\vert +\alpha_{2}^{2}\left\vert \left(
\begin{array}
[c]{cc}%
w_{21} & w_{22}%
\end{array}
\right)  \right\rangle \left\langle \left(
\begin{array}
[c]{cc}%
w_{21}^{\boldsymbol{A}} & w_{22}^{\boldsymbol{A}}%
\end{array}
\right)  \right\vert
\end{align}%
\begin{align}
&  \boldsymbol{d}g\left(  \boldsymbol{r}\right)  ^{\dagger}\boldsymbol{d}%
g\left(  \boldsymbol{r}\right)  =\left(
\begin{array}
[c]{ccc}%
2x & 2y & 2z\\
2x & 2y & 0
\end{array}
\right)  ^{\dagger}\left(
\begin{array}
[c]{ccc}%
2x & 2y & 2z\\
2x & 2y & 0
\end{array}
\right) \nonumber\\
&  =\left(
\begin{array}
[c]{cc}%
2x & 2x\\
2y & 2y\\
2z & 0
\end{array}
\right)  \left(
\begin{array}
[c]{ccc}%
2x & 2y & 2z\\
2x & 2y & 0
\end{array}
\right) \\
&  =\left(
\begin{array}
[c]{ccc}%
16x^{2} & 16xy & 4xz\\
16xy & 16y^{2} & 4yz\\
4xz & 4yz & 0
\end{array}
\right) \nonumber\\
&  =\alpha_{1}^{2}\left\vert \left(
\begin{array}
[c]{ccc}%
v_{11} & v_{12} & v_{13}%
\end{array}
\right)  \right\rangle \left\langle \left(
\begin{array}
[c]{ccc}%
v_{11}^{\boldsymbol{A}} & v_{12}^{\boldsymbol{A}} & v_{13}^{\boldsymbol{A}}%
\end{array}
\right)  \right\vert \\
&  +\alpha_{2}^{2}\left\vert \left(
\begin{array}
[c]{ccc}%
v_{21} & v_{22} & v_{23}%
\end{array}
\right)  \right\rangle \left\langle \left(
\begin{array}
[c]{ccc}%
v_{21}^{\boldsymbol{A}} & v_{22}^{\boldsymbol{A}} & v_{23}^{\boldsymbol{A}}%
\end{array}
\right)  \right\vert
\end{align}%
\begin{align}
&  \mathcal{N}_{\left(  x,y,z\right)  }\boldsymbol{g\quad\bot\quad}%
\mathcal{T}_{\left(  x,y,z\right)  }\boldsymbol{g}\\
&  \mathcal{T}_{\left(  x,y,z\right)  }\boldsymbol{g}=\left[  \mathcal{N}%
_{\left(  x,y,z\right)  }\boldsymbol{g}\right]  ^{\bot}%
\end{align}%
\begin{equation}
\operatorname*{Rank}\left\{  d\boldsymbol{g}\left(  \boldsymbol{x}\right)
\right\}  =\dim\left\{  \mathcal{N}_{\left(  x,y,z\right)  }\boldsymbol{g}%
\right\}
\end{equation}
e.g. 4%

\begin{align}
g  &  :\mathbb{R}^{3}\rightarrow\mathbb{R}^{2}\nonumber\\
g_{1}\left(  x,y,z\right)   &  =x^{2}+y^{2}+z^{2}=c_{1}^{2}\nonumber\\
g_{2}\left(  x,y,z\right)   &  =x^{2}-y^{2}=c_{2}^{2}\nonumber\\
\left\langle \nabla g_{1}\left(  x,y,z\right)  \right\vert  &  =\left(
2x,2y,2z\right) \nonumber\\
\left\langle \nabla g_{2}\left(  x,y,z\right)  \right\vert  &  =\left(
2x,-2y,0\right) \nonumber\\
\ker\left\{  d\boldsymbol{g}\left(  x,y,z\right)  \right\}   &  =\left\{
\left.  \left(
\begin{array}
[c]{c}%
\delta x\\
\delta y\\
\delta z
\end{array}
\right)  \right\vert \left(
\begin{array}
[c]{ccc}%
2x & 2y & 2z\\
2x & -2y & 0
\end{array}
\right)  \left(
\begin{array}
[c]{c}%
\delta x\\
\delta y\\
\delta z
\end{array}
\right)  =\left(
\begin{array}
[c]{c}%
0\\
0
\end{array}
\right)  \right\} \nonumber\\
\mathcal{N}_{\left(  x,y,z\right)  }\boldsymbol{g}  &
=\operatorname*{linearspan}\left\{  \left(  2x,2y,2z\right)  ,\left(
2x,-2y,0\right)  \right\}
\end{align}
e.g.5%

\begin{align}
g_{1}  &  :\mathbb{R}^{2}\rightarrow\mathbb{R}\nonumber\\
g_{1}\left(  x,y\right)   &  =\left(  x^{2}+y^{2}-1\right)  ^{2}=0=x^{4}%
+y^{4}+4x^{2}y^{2}-2x^{2}-2y^{2}+1\nonumber\\
\left\langle \nabla g_{1}\left(  x,y\right)  \right\vert  &  =4\left(  \left(
x^{2}+y^{2}-1\right)  x,\left(  x^{2}+y^{2}-1\right)  y\right) \nonumber\\
\ker\left\{  \boldsymbol{d}g_{1}\left[  \boldsymbol{r}_{0}\right]  \right\}
&  =\left\{
\begin{array}
[c]{c}%
\left.  \left(
\begin{array}
[c]{c}%
\delta x\\
\delta y
\end{array}
\right)  \right\vert \left(
\begin{array}
[c]{cc}%
\left(  x^{2}+y^{2}-1\right)  x & \left(  x^{2}+y^{2}-1\right)  y
\end{array}
\right)  \left(
\begin{array}
[c]{c}%
\delta x\\
\delta y
\end{array}
\right) \\
=\left(
\begin{array}
[c]{c}%
0\\
0
\end{array}
\right)
\end{array}
\right\} \nonumber\\
\mathcal{N}_{\left(  x,y\right)  }g_{1}  &  =\left[  \mathcal{T}_{\left(
x,y\right)  }g_{1}\right]  ^{\bot}%
\end{align}
Let $\left(  x,y\right)  =\left(  0,1\right)  $ then
\begin{equation}
\left(
\begin{array}
[c]{cc}%
\left(  x^{2}+y^{2}-1\right)  x & \left(  x^{2}+y^{2}-1\right)  y
\end{array}
\right)  =\left(
\begin{array}
[c]{cc}%
0 & 0
\end{array}
\right)
\end{equation}
which gives
\begin{equation}
\ker\left\{  dg_{1}\left[  \left(  0,1\right)  \right]  \right\}
=\mathbb{R}^{2}.
\end{equation}
A tangent vector $\left(
\begin{array}
[c]{c}%
\delta x\\
\delta y
\end{array}
\right)  $ at the point $\boldsymbol{r}_{0}$ is defined by considering curves%
\begin{align}
g_{1}\left(  v\left(  t\right)  \right)   &  =0\,,t\in\left[  0,\delta\right]
\nonumber\\
v\left(  0\right)   &  =\boldsymbol{r}_{0},
\end{align}
$%
\begin{array}
[c]{c}%
\text{such curves could be }x\left(  t\right)  =\cos\left(  \omega
t+\omega_{0}\right)  ;y\left(  t\right)  =\sin\left(  \omega t+\omega
_{0}\right)  ;t\in\left[  0,\frac{2\pi}{\omega}\right)  \text{, }\\
\omega\text{ is an arbitary constant }\\
\text{and }\omega_{0}\text{ determines the point where the tangent space is
evaluated,}%
\end{array}
$

and must be a solution of
\begin{equation}
\left(
\begin{array}
[c]{c}%
\delta x\\
\delta y
\end{array}
\right)  =\left(
\begin{array}
[c]{c}%
\frac{\boldsymbol{d}v_{x}\left(  t\right)  }{\boldsymbol{d}t}\\
\frac{\boldsymbol{d}v_{y}\left(  t\right)  }{\boldsymbol{d}t}%
\end{array}
\right)
\end{equation}
This implies the necessary condition
\begin{equation}
\frac{dg_{1}\left(  v\left(  t\right)  \right)  }{dt}=\left\langle \nabla
g_{1}\left(  x,y\right)  \left\vert \frac{dv\left(  t\right)  }{dt}\right.
\right\rangle =\nabla g_{1}\left(  x,y\right)  ^{t}\left(
\begin{array}
[c]{c}%
\frac{dx}{dt}\\
\frac{dy}{\boldsymbol{d}t}%
\end{array}
\right)  =0
\end{equation}
i.e.$\left(
\begin{array}
[c]{c}%
\delta x\\
\delta y
\end{array}
\right)  \in$ $\ker\left\{  dg_{1}\left[  \boldsymbol{r}_{0}\right]  \right\}
$. The curve $g_{1}$ is the unit circles with center at the origin, thus its
tangent at $\left(  0,1\right)  $ is the line generated by $\left(
1,0\right)  $ i.e the vectors $\left\{  \left(
\begin{array}
[c]{c}%
\delta x\\
0
\end{array}
\right)  \right\}  ,$which form a subspace of $\ker\left\{  dg_{1}\left[
\left(  0,1\right)  \right]  \right\}  .$ This can be seen by considering the
curve
\begin{equation}
x\left(  t\right)  =\cos\left(  \omega t+\omega_{0}\right)  ,\,y\left(
t\right)  =\sin\left(  \omega t+\omega_{0}\right)  ,
\end{equation}
where
\begin{align}
x  &  =x\left(  0\right)  =\cos\left(  \omega_{0}\right)  \Rightarrow
\omega_{0}=\cos^{-1}\left(  x\right) \nonumber\\
y  &  =y\left(  0\right)  =\sin\left(  \omega_{0}\right)  \Rightarrow
\omega_{0}=\sin^{-1}\left(  y\right)
\end{align}
which has derivatives
\begin{equation}
\left(
\begin{array}
[c]{c}%
\frac{\boldsymbol{d}x}{\boldsymbol{d}t}\\
\frac{\boldsymbol{d}y}{\boldsymbol{d}t}%
\end{array}
\right)  =\left(
\begin{array}
[c]{c}%
-\omega\sin\left(  \omega t+\omega_{0}\right) \\
\omega\cos\left(  \omega t+\omega_{0}\right)
\end{array}
\right)
\end{equation}
and thus tangent vectors at the point $\left(  x,y\right)  =\left(  x\left(
0\right)  ,y\left(  0\right)  \right)  $
\begin{equation}
\left.  \left(
\begin{array}
[c]{c}%
\frac{\boldsymbol{d}x}{\boldsymbol{d}t}\\
\frac{\boldsymbol{d}y}{\boldsymbol{d}t}%
\end{array}
\right)  \right\vert _{t=0}=\left(
\begin{array}
[c]{c}%
-\omega\sin\left(  \omega_{0}\right) \\
\omega\cos\left(  \omega_{0}\right)
\end{array}
\right)  =\omega\left(
\begin{array}
[c]{c}%
-y\\
x
\end{array}
\right)  .
\end{equation}
At the point $\left(  0,1\right)  \equiv\omega_{0}=\frac{\pi}{2}$ this gives
\begin{equation}
\left.  \left(
\begin{array}
[c]{c}%
\frac{\boldsymbol{d}x}{\boldsymbol{d}t}\\
\frac{\boldsymbol{d}y}{\boldsymbol{d}t}%
\end{array}
\right)  \right\vert _{t=0}=\left(
\begin{array}
[c]{c}%
-\omega\\
0
\end{array}
\right)  ;\omega\in\mathbb{R}%
\end{equation}
and thus
\begin{equation}
\mathcal{T}_{\left(  0,1\right)  }g_{1}=\operatorname*{linearspan}\left\{
\left(
\begin{array}
[c]{cc}%
1 & 0
\end{array}
\right)  \right\}  .
\end{equation}
$dg_{1}\left[  \left(  0,1\right)  \right]  $ $\equiv0$ is not surjective and
does not equal the normal space which is defined as
\begin{equation}
\mathcal{N}_{\left(  x,y\right)  }g_{1}=\left[  \mathcal{T}_{\left(
x,y\right)  }g_{1}\right]  ^{\bot}=\operatorname*{linearspan}\left\{  \left(
\begin{array}
[c]{cc}%
0 & 1
\end{array}
\right)  \right\}  .
\end{equation}


e.g.6%

\begin{align}
g_{1}  &  :\mathbb{R}^{2}\rightarrow\mathbb{R}\nonumber\\
&  g_{1}\left(  x,y\right)  =x^{2}-y^{2}=0\nonumber\\
&  \left\langle \nabla g_{1}\left(  x,y\right)  \right\vert =2\left(
x,-y\right) \nonumber\\
&  \ker\left\{  dg_{1}\left[  \boldsymbol{r}_{0}\right]  \right\}  =\left\{
\left.  \left(
\begin{array}
[c]{c}%
\delta x\\
\delta y
\end{array}
\right)  \right\vert \left(
\begin{array}
[c]{cc}%
x & -y
\end{array}
\right)  \left(
\begin{array}
[c]{c}%
\delta x\\
\delta y
\end{array}
\right)  =\left(
\begin{array}
[c]{c}%
0\\
0
\end{array}
\right)  \right\} \nonumber\\
&  \mathcal{N}_{\left(  x,y\right)  }g_{1}=\operatorname*{linearspan}\left\{
\left(
\begin{array}
[c]{cc}%
x & -y
\end{array}
\right)  \right\}
\end{align}
Let $\left(  x,y\right)  =\left(  0,0\right)  $ then
\begin{equation}
\nabla g_{1}\left(  0,0\right)  =\left(
\begin{array}
[c]{cc}%
x & -y
\end{array}
\right)  =\left(
\begin{array}
[c]{cc}%
0 & 0
\end{array}
\right)
\end{equation}
which also gives
\begin{equation}
\ker\left\{  dg_{1}\left[  \left(  0,1\right)  \right]  \right\}
=\mathbb{R}^{2}.
\end{equation}
Now
\begin{align}
\mathcal{T}_{\left(  0,0\right)  }g_{1}  &  =\operatorname*{linearspan}%
\left\{  \left(
\begin{array}
[c]{cc}%
0 & 1
\end{array}
\right)  \right\} \nonumber\\
\mathcal{N}_{\left(  x,y\right)  }g_{1}  &  =\left[  \mathcal{T}_{\left(
0,0\right)  }g_{1}\right]  ^{\bot}=\operatorname*{linearspan}\left\{  \left(
\begin{array}
[c]{cc}%
1 & 0
\end{array}
\right)  \right\}
\end{align}


A point on the constraint manifold is called \emph{normal }if $\left\{  \nabla
g_{1}\left(  x,y\right)  \right\}  $ generates the normal space, which at the
point $\left(  0,0\right)  $ it does not, and \emph{regular }if $\mathcal{T}%
_{\left(  x,y\right)  }g_{1}=\ker\left\{  dg_{1}\left[  \left(  x,y\right)
\right]  \right\}  ,$ which again at the point $\left(  0,0\right)  $ is not
true. \emph{Normal implies regular. }Thus for a normal point $\left(
x,y\right)  $
\begin{align}
\mathcal{N}_{\left(  x,y\right)  }g_{1}  &  =\left[  \mathcal{T}_{\left(
x,y\right)  }g_{1}\right]  ^{\bot}=\ker\left\{  dg_{1}\left[  \left(
x,y\right)  \right]  \right\}  ^{\bot}\nonumber\\
\ker\left\{  dg_{1}\left[  \left(  x,y\right)  \right]  \right\}  ^{\bot}  &
\cong\operatorname*{range}\left\{  dg_{1}\left[  \left(  x,y\right)  \right]
\right\}
\end{align}


\begin{itemize}
\item $m=n$
\end{itemize}

\begin{itemize}
\item $m<n$

\item $m>n$
\end{itemize}

\section{Banach Space}

If a Banach space $\mathcal{B}$ has a basis $\left\{  \left.  \left\vert
e_{j}\right\rangle \right\vert 1\leq j\leq\infty\right\}  $ then
\[
T:\mathcal{B}\rightarrow\mathbb{R}^{\infty}%
\]
where $\mathbb{R}^{\infty}$ is the space of sequencies (not in general
convergent)\ given by
\[
\left\vert c\right\rangle =\sum_{1\leq j\leq\infty}c_{j}\left\vert
e_{j}\right\rangle \equiv\left\vert \boldsymbol{e}\right\rangle \boldsymbol{c}%
\]
this correspondence might not be $1-1.$ If the Banach space is a closure of a
subset of distributions then
\[
\left\vert c\right\rangle =\int c\left(  x\right)  \left\vert x\right\rangle
\equiv\left\vert \boldsymbol{x}\right\rangle \boldsymbol{c}%
\]


\section{\bigskip Extra}

As the example is quadratic it useful to consider the Hessian $d^{2}%
\boldsymbol{f}\left(  A\right)  $ and note that a sufficient condition that
$\boldsymbol{A}+\delta\boldsymbol{A}$ $\in\mathfrak{M}$ is that $\delta
A\in\ker\left\{  d\boldsymbol{g}\left(  \boldsymbol{A}\right)  \right\}  $
and
\begin{equation}
\left(  \left.  \delta A\right\vert \widehat{\Phi^{\dagger}\Phi_{\varphi}%
}\delta A\right)  =\mathbf{0}.
\end{equation}
$\lceil$%
\begin{equation}
\left(  \left.  \delta A\right\vert \Phi\left(  \varphi_{j}\right)  ^{\dagger
}\Phi\left(  \varphi_{j}\right)  \delta A\right)  =\operatorname{Tr}\left\{
\Phi\left(  \varphi_{j}\right)  ^{\dagger}\Phi\left(  \varphi_{j}\right)
\left(  \delta A\right)  ^{2}\right\}  =0
\end{equation}
In order to satisfy this constraint a positive matrix to have a zero diagonal,
which is not possible as the trace would then be zero%
\begin{align}
\mathbf{\rho}  &  =\operatorname{Tr}\left\{  \Phi^{\dagger}\Phi_{\varphi}%
D_{1}\right\}  =\operatorname{Tr}\left\{  \Phi^{\dagger}\Phi_{\varphi}%
D_{2}\right\} \\
\Delta &  =D_{1}-D_{2}\\
\mathbf{0}  &  =\operatorname{Tr}\left\{  \Phi^{\dagger}\Phi_{\varphi}%
\Delta\right\}
\end{align}
$\rfloor$

Misc:
\begin{align*}
\mathbb{R}^{m}  &  =\operatorname*{Range}\left\{  d\boldsymbol{g}\left(
\boldsymbol{A}\right)  \right\}  \oplus\operatorname*{Range}\left\{
d\boldsymbol{g}\left(  \boldsymbol{A}\right)  \right\}  ^{c}\\
&  \qquad\cup\qquad\qquad\qquad\qquad\cap\\
\mathbb{R}^{m}  &  =\left[  \ker\left\{  d\boldsymbol{g}\left(  \boldsymbol{A}%
\right)  \right\}  ^{c}\right]  \oplus\left[  \ker\left\{  d\boldsymbol{g}%
\left(  \boldsymbol{A}\right)  \right\}  ^{c}\right]  ^{c},
\end{align*}
which gives for inner product (self dual) spaces
\[
\mathbb{R}^{m}=\ker\left\{  d\boldsymbol{g}\left(  \boldsymbol{A}\right)
\right\}  ^{\bot}\oplus\ker\left\{  d\boldsymbol{g}\left(  \boldsymbol{A}%
\right)  \right\}  .
\]



\end{document}